% Options for packages loaded elsewhere
\PassOptionsToPackage{unicode}{hyperref}
\PassOptionsToPackage{hyphens}{url}
\documentclass[
  a4paper,
]{ltjsarticle}
\usepackage{xcolor}
\usepackage[margin=25mm]{geometry}
\usepackage{amsmath,amssymb}
\setcounter{secnumdepth}{-\maxdimen} % remove section numbering
\usepackage{iftex}
\ifPDFTeX
  \usepackage[T1]{fontenc}
  \usepackage[utf8]{inputenc}
  \usepackage{textcomp} % provide euro and other symbols
\else % if luatex or xetex
  \usepackage{unicode-math} % this also loads fontspec
  \defaultfontfeatures{Scale=MatchLowercase}
  \defaultfontfeatures[\rmfamily]{Ligatures=TeX,Scale=1}
\fi
\usepackage{lmodern}
\ifPDFTeX\else
  % xetex/luatex font selection
\fi
% Use upquote if available, for straight quotes in verbatim environments
\IfFileExists{upquote.sty}{\usepackage{upquote}}{}
\IfFileExists{microtype.sty}{% use microtype if available
  \usepackage[]{microtype}
  \UseMicrotypeSet[protrusion]{basicmath} % disable protrusion for tt fonts
}{}
\makeatletter
\@ifundefined{KOMAClassName}{% if non-KOMA class
  \IfFileExists{parskip.sty}{%
    \usepackage{parskip}
  }{% else
    \setlength{\parindent}{0pt}
    \setlength{\parskip}{6pt plus 2pt minus 1pt}}
}{% if KOMA class
  \KOMAoptions{parskip=half}}
\makeatother
\usepackage{longtable,booktabs,array}
\newcounter{none} % for unnumbered tables
\usepackage{calc} % for calculating minipage widths
% Correct order of tables after \paragraph or \subparagraph
\usepackage{etoolbox}
\makeatletter
\patchcmd\longtable{\par}{\if@noskipsec\mbox{}\fi\par}{}{}
\makeatother
% Allow footnotes in longtable head/foot
\IfFileExists{footnotehyper.sty}{\usepackage{footnotehyper}}{\usepackage{footnote}}
\makesavenoteenv{longtable}
\usepackage{graphicx}
\makeatletter
\newsavebox\pandoc@box
\newcommand*\pandocbounded[1]{% scales image to fit in text height/width
  \sbox\pandoc@box{#1}%
  \Gscale@div\@tempa{\textheight}{\dimexpr\ht\pandoc@box+\dp\pandoc@box\relax}%
  \Gscale@div\@tempb{\linewidth}{\wd\pandoc@box}%
  \ifdim\@tempb\p@<\@tempa\p@\let\@tempa\@tempb\fi% select the smaller of both
  \ifdim\@tempa\p@<\p@\scalebox{\@tempa}{\usebox\pandoc@box}%
  \else\usebox{\pandoc@box}%
  \fi%
}
% Set default figure placement to htbp
\def\fps@figure{htbp}
\makeatother
\setlength{\emergencystretch}{3em} % prevent overfull lines
\providecommand{\tightlist}{%
  \setlength{\itemsep}{0pt}\setlength{\parskip}{0pt}}
% Glyph map for lualatex (newunicodechar). Maps symbols that may lack font glyphs.
\usepackage{newunicodechar}
\usepackage{amssymb}
\newunicodechar{∎}{\ensuremath{\blacksquare}}
\newunicodechar{✓}{\ensuremath{\checkmark}}
\newunicodechar{✗}{\ensuremath{\times}}
\newunicodechar{⟹}{\ensuremath{\Longrightarrow}}
\newunicodechar{⟺}{\ensuremath{\Longleftrightarrow}}
\newunicodechar{↪}{\ensuremath{\hookrightarrow}}
\newunicodechar{⌊}{\ensuremath{\lfloor}}
\newunicodechar{⌋}{\ensuremath{\rfloor}}
\newunicodechar{≃}{\ensuremath{\simeq}}
\newunicodechar{≈}{\ensuremath{\approx}}
\newunicodechar{≤}{\ensuremath{\leq}}
\newunicodechar{≥}{\ensuremath{\geq}}
\newunicodechar{≠}{\ensuremath{\neq}}
\newunicodechar{→}{\ensuremath{\rightarrow}}
\newunicodechar{←}{\ensuremath{\leftarrow}}
\newunicodechar{↔}{\ensuremath{\leftrightarrow}}
\newunicodechar{⊥}{\ensuremath{\perp}}
\newunicodechar{√}{\ensuremath{\surd}}
\newunicodechar{∑}{\ensuremath{\sum}}
\newunicodechar{∏}{\ensuremath{\prod}}
\newunicodechar{∞}{\ensuremath{\infty}}
\newunicodechar{∈}{\ensuremath{\in}}
\newunicodechar{∀}{\ensuremath{\forall}}
\newunicodechar{∣}{\ensuremath{\mid}}
\newunicodechar{γ}{\ensuremath{\gamma}}
\newunicodechar{λ}{\ensuremath{\lambda}}
\newunicodechar{μ}{\ensuremath{\mu}}
\newunicodechar{ρ}{\ensuremath{\rho}}
\newunicodechar{σ}{\ensuremath{\sigma}}
\newunicodechar{φ}{\ensuremath{\varphi}}
\newunicodechar{ψ}{\ensuremath{\psi}}
\newunicodechar{θ}{\ensuremath{\theta}}
\newunicodechar{Δ}{\ensuremath{\Delta}}
\newunicodechar{Π}{\ensuremath{\Pi}}
\newunicodechar{Λ}{\ensuremath{\Lambda}}
\newunicodechar{τ}{\ensuremath{\tau}}
\newunicodechar{ε}{\ensuremath{\varepsilon}}
\newunicodechar{ω}{\ensuremath{\omega}}
\newunicodechar{π}{\ensuremath{\pi}}
\newunicodechar{ν}{\ensuremath{\nu}}
\newunicodechar{±}{\ensuremath{\pm}}
\newunicodechar{×}{\ensuremath{\times}}
\newunicodechar{·}{\ensuremath{\cdot}}
\usepackage{bookmark}
\IfFileExists{xurl.sty}{\usepackage{xurl}}{} % add URL line breaks if available
\urlstyle{same}
\hypersetup{
  pdftitle={Dynamics of Self-Consistent Relational-Wave Closed Systems------From the local instability of the 90-degree relative equilibrium to inflation{,} spontaneous symmetry breaking{,} Grassmann geometry{,} terminal universality{,} and relationalism},
  pdfauthor={Noriaki Kihara (WF System Co., Ltd.), ORCID 0009-0004-6753-4020},
  hidelinks,
  pdfcreator={LaTeX via pandoc}}

\title{Dynamics of Self-Consistent Relational-Wave Closed
Systems------From the local instability of the 90-degree relative
equilibrium to inflation, spontaneous symmetry breaking, Grassmann
geometry, terminal universality, and relationalism}
\author{Noriaki Kihara (WF System Co., Ltd.), ORCID 0009-0004-6753-4020}
\date{2026-09-12}

\begin{document}
\maketitle

\textbf{Series:} Foundations of Self-Consistent Relational-Wave Closed
Systems (Paper 0, capstone / of 10)\\
\textbf{Author:} Noriaki Kihara (WF System Co., Ltd.)　\textbf{Date:}
2026-09-12\\
\textbf{Version DOI:} 10.5281/zenodo.22729147\\
\textbf{Concept DOI:} 10.5281/zenodo.22729146\\
\textbf{ORCID:} 0009-0004-6753-4020\\
\textbf{Zenodo:} https://zenodo.org/records/22729147\\

\begin{center}\rule{0.5\linewidth}{0.5pt}\end{center}

\subsection{Abstract}\label{abstract}

For the dynamics defined in the prior work (Mechanism of inflationary
rapid expansion in self-consistent relational-wave closed systems,
Concept DOI 10.5281/zenodo.22112008)------place a complex relational
wave \(z_e\) on each edge of the complete graph \(K_N\), and drive it
with the real orthogonal map \(z\to\exp(\Delta\tau K)z\)
(\(\Delta\tau=2\pi/N\)) generated by the phase-only real antisymmetric
generator \(K_{ef}=A_{ef}\sin(\varphi_f-\varphi_e)\) (\(A\) is the
adjacency of the line graph \(L(K_N)\))------this capstone establishes
the nine foundational papers (Papers 1--9) as \textbf{a single causal
chain}:

\[\text{map}\ \to\ \text{90-degree self-consistent relative equilibrium (floor)}\ \to\ \text{ignition qualification = local instability of the floor}\ \to\ \textbf{inflationary rapid expansion}\ \to\ \text{SSB / plane reorientation}\ \to\ \text{terminal amplitude universalization}\ \to\ \text{only relational invariants remain}.\]

\textbf{Central proposition}: the local instability of the
self-consistent 90-degree relative equilibrium produces
\textbf{inflationary rapid expansion} (exponential growth of the
out-of-floor-plane component \(H_\perp/H\)), spontaneous symmetry
breaking, and plane reorientation; after nonlinear saturation it
universalizes each wave's amplitude to \(\sqrt{H/M}\) while leaving only
the relational phase/geometric information. The \textbf{mechanism} of
the rapid expansion is the linear instability of the floor, and its
\textbf{real amplification rate} is refined by the finite-time Lyapunov
rate of the time-varying tangent cocycle along the real orbit
(phenomenon = inflation, mechanism = linear instability + time-varying
tangent cocycle).

Each stage of the chain: (1) exact derivation of the single-wave
equation of motion (Kuramoto-type phase coupling; amplitude driven by
the second harmonic \(\sin2\Delta\varphi\); Paper 1). (2) On the
90-degree skeleton, \(D^{-1}KD=iW\)
(\(D=\operatorname{diag}(e^{i\varphi_e})\), \(W\) real symmetric
nonnegative), so the complex self-consistent floor problem is exactly
equivalent to the Perron problem of the real symmetric \(W\); a
positive-amplitude floor exists uniquely, and from the parity bipartite
structure \(\sigma^2\), each edge amplitude², and the full spectrum can
be written in closed form for \textbf{arbitrary \(N\)}, with active
dimension \(O(N)\) (Papers 1 and 3). (3) The ignition qualification is
``an unstable relative equilibrium'' only; high symmetry, zero centroid,
and square-closure are not qualifications (Paper 2). (4) \textbf{On the
cyclic-Fourier high-symmetry floor} the factorization of zero-closure is
quantized into discrete block sizes with \(L=N/\gcd(N,4)\),
\(p=\operatorname{spf}(L)\), and \(p=2\iff 8\mid N\) (\(\nu_2(N)\ge3\))
is the 2-adic selection rule of the square map (Paper 4). (5) An
infinitesimal seed breaks the symmetry and triggers \textbf{inflationary
rapid expansion}. It begins with the floor-Jacobian initial/frozen rate
\(\lambda_0=\ln\rho_0\) (\(\rho_0=\rho(J_0)\), \(=\) the Landau summit
curvature \(V''(0)=-\lambda_0\)), and its \textbf{real amplification
rate is the finite-time Lyapunov rate \(\gamma_t\) of the tangent
cocycle \(\Phi(t,0)=\prod_t J(z_t)\) along the real orbit} (for
\(N=6,8,10\), \(\gamma_t>\lambda_0\)) (Papers 5 and 6). (6) The state is
a rank-2 plane \(P(t)\) at each instant
(\(\Pi_{t+1}=Q_t\Pi_tQ_t^{\mathsf T}\)), with
\(H_\perp/H=\tfrac14\|\Pi_t-\Pi_0\|_F^2\)------inflation is exponential
departure from the initial plane on \(\mathrm{Gr}(2,M)\), and the
terminal plane spans four dimensions with the original plane via two
nonzero principal angles (the cumulative map is not \(SO(4)\); Paper 7).
(7) The terminus universalizes to equal amplitude
\(a_\infty=\sqrt{H/M}\) (a state quantity, extensive \((H,M)\) /
intensive \(\rho=H/M\)) and \(\Delta_{\rm tail}=-2\pi/N\) (Papers 5 and
8). The \textbf{degenerate vacuum family} is not identifiable as a
finite discrete orbit of \(S_N\) + global phase + conjugation (it was
confirmed that the five terminal states for \(N=6\) do not map to one
another); the orientation is seed-dependent (SSB confirmed in all 35
values of \(N\) × 5 seeds = 175 runs). That it is not a finite discrete
orbit is established; continuity and the identification of a larger
continuous symmetry group are open (Paper 5). (8) The waves are
anonymous (\(S_N\) gauge), and closure alone does not generate discrete
species (\(\mathcal C_k\simeq Q_{k-2}/S_k\), continuous for \(k\ge3\));
what is discrete is the block size \(p\) that the floor quantizes, and
the physical content resides in relational invariants (Paper 9).

These are organized along three unifying axes: \textbf{A. Unification of
instability}------ignition qualification, SSB summit curvature,
inflationary amplification, and Grassmann departure all begin from a
single local instability \(\rho(J_0)>1\), and the tangent cocycle
carries the subsequent real amplification. \textbf{B. Memory and
erasure}------the transient records the initial structure in
\(J_0,\Phi(t,0)\), but the terminal amplitude leaves only \(H/M\) and
forgets the configuration, while the SSB history remains in the phase
configuration and plane orientation (selective erasure of history
information). \textbf{C. Relationalism}------neither an absolute wave
number nor an absolute normal can be read; the physical content resides
only in relational invariants modulo \(S_N\), global phase, and
coordinate rotation.

Note that the terms ``inflation,'' ``vacuum,'' and ``spontaneous
symmetry breaking'' used in this paper are designations for mathematical
features of this dynamical system; a substantive correspondence with
cosmology or particle physics is open, as stated in §9.

This capstone stands on nine foundational papers (Papers 1--9). Its
scope is limited to self-consistent single-coherent relative equilibria
with a 90-degree skeleton and their small anisotropic perturbations; it
does not enter into the identification of coordinates (whether \(xyzt\)
or including \(R,Q\)), the type of gauge (equal-interval / logarithmic /
winding), the multi-particle case, or the understanding of the vacuum
itself. These are explicitly noted as open (§9).

\begin{center}\rule{0.5\linewidth}{0.5pt}\end{center}

\subsection{1. Introduction}\label{introduction}

This research series has numerically studied the dynamics of a closed
system in which all pairwise relations \(M=N(N-1)/2\) among \(N\)
entities are taken to be complex relational waves \(z_e\in\mathbb C\),
with external seeds removed. The waves live on the edges of the complete
graph \(K_N\), and the interaction is determined only by the phase
differences between edges that share a vertex. The prior paper {[}1{]}
reported this system's inflationary rapid expansion (exponential growth
of the out-of-parent-plane component \(H_\perp/H\), rank-4-like readout,
transition to a metastable structure), but did not show \textbf{why a
single relational wave takes that orbit} at the level of the equation of
motion. Taking that single point as its starting point, this capstone
establishes as a single causal chain the necessity of the initial value
(the 90-degree floor), its qualification and exact solution, the reading
of the rapid expansion as spontaneous symmetry breaking, the Grassmann
geometry after locking, and the terminal amplitude universalization and
relationalism.

The backbone of this capstone is \textbf{a single causal chain} running
through Papers 1--9 (Papers 4, 8, and 9 are not auxiliary but the body
of the chain).

\begin{itemize}
\tightlist
\item
  \textbf{(a) Dynamics} (§2, Paper 1): each edge wave couples in
  Kuramoto fashion with the vertex-sharing adjacent edges; the phase
  advances monotonically as \(\sin^2\), and the amplitude increases or
  decreases via the second harmonic \(\sin(2\Delta\varphi)\).
\item
  \textbf{(b) Floor = 90-degree self-consistent relative equilibrium and
  its exact solution} (§3, §4, Papers 1 and 3): the maximally symmetric
  relative equilibrium that freezes the amplitude termwise is limited to
  the 90-degree lattice (the analytic reason why the termwise floor
  reached by make\_parent becomes a 90-degree lattice; why the iteration
  selects that termwise branch is open, §4). On the 90-degree skeleton,
  \(D^{-1}KD=iW\) makes the floor exactly equivalent to the Perron
  problem of the real symmetric \(W\); from the parity bipartite
  structure \(\sigma^2\), amplitude², and the full spectrum are closed
  form for \textbf{arbitrary \(N\)}, and the active dimension is
  \(O(N)\).
\item
  \textbf{(c) Ignition qualification and quantization of zero-closure}
  (§3.4, §3.5, Papers 2 and 4): the ignition qualification is ``an
  unstable relative equilibrium'' only (high symmetry, zero centroid,
  and square-closure are not qualifications). Zero-closure
  \(\sum z^2=0\) alone is a continuous moduli, but the discrete symmetry
  of the cyclic-Fourier high-symmetry floor quantizes it into a block
  size \(p=\operatorname{spf}(N/\gcd(N,4))\) (\(p=2\iff 8\mid N\)).
\item
  \textbf{(d) Inflationary rapid expansion------from local instability
  to the tangent cocycle} (§5, §6, Papers 5 and 6): an infinitesimal
  seed breaks the symmetry, and the out-of-floor-plane component
  \(H_\perp/H\) grows exponentially (\textbf{inflation}). The rapid
  expansion begins with the floor-Jacobian initial/frozen rate
  \(\lambda_0=\ln\rho_0\) (\(\rho_0=\rho(J_0)\), \(=\) the Landau summit
  curvature \(V''(0)=-\lambda_0\)), and its \textbf{real amplification
  rate is promoted to the finite-time Lyapunov rate \(\gamma_t\) of the
  tangent cocycle \(\Phi(t,0)=\prod_t J(z_t)\) along the real orbit}
  (for \(N=6,8,10\), \(\gamma_t>\lambda_0\)).
\item
  \textbf{(e) Geometry------exponential departure on the Grassmann} (§7,
  Paper 7): the state is a rank-2 plane \(P(t)\) at each instant
  (\(\Pi_{t+1}=Q_t\Pi_tQ_t^{\mathsf T}\)), and since
  \(H_\perp/H=\tfrac14\|\Pi_t-\Pi_0\|_F^2\), inflation is exponential
  departure from the initial plane on \(\mathrm{Gr}(2,M)\); the terminal
  plane spans four dimensions with two nonzero principal angles (the
  cumulative map is not \(SO(4)\)).
\item
  \textbf{(f) Terminal universalization and relationalism} (§8, Papers
  5, 8, and 9): the terminus universalizes to equal amplitude
  \(a_\infty=\sqrt{H/M}\) (an intensive terminal state quantity) and
  forgets the configuration, while the SSB history remains in the phase
  configuration and plane orientation (selective erasure of memory). The
  waves are anonymous (\(S_N\) gauge), and the physical content resides
  only in relational invariants.
\end{itemize}

These (a)--(f) converge onto three unifying axes: \textbf{A. Unification
of instability} (ignition qualification, SSB summit curvature,
inflationary amplification, and Grassmann departure all begin from a
single local instability \(\rho(J_0)>1\), and the tangent cocycle
carries the real amplification), \textbf{B. Memory and erasure} (the
transient records the initial structure, the terminal amplitude forgets,
and the phase/geometry retains the history), \textbf{C. Relationalism}
(neither an absolute wave number nor an absolute normal can be read;
only relational invariants are physical).

What is introduced and used in this paper is only the dynamics defined
in the prior paper (§2). Zero-closure and a particular amplitude
distribution are not assumed but derived. Each section carries its
claims and the key formulas/figures within this capstone, and delegates
detailed derivations, exhaustive data, additional figures, and
reproduction procedures to the corresponding foundational papers (Papers
1--9; Concept DOIs in Appendix A). The ``density'' of this capstone lies
in integrating the nine papers as a single chain, not in duplicating raw
data.

\subsection{2. Dynamics and the single-wave equation of
motion}\label{dynamics-and-the-single-wave-equation-of-motion}

\subsubsection{2.1 The map}\label{the-map}

The one\_step of the prior paper is given as follows. Let \(A\) be the
adjacency matrix of the line graph \(L(K_N)\) of \(K_N\) (\(A_{ef}=1\)
if edges \(e,f\) share a vertex), \(\varphi_e=\arg z_e\), then
\[K_{ef}=A_{ef}\sin(\varphi_f-\varphi_e)\] is a real antisymmetric
matrix (its entries are only the sines of phase differences, not
containing the amplitude \(|z_e|\)). one\_step is
\[z_{n+1}=\exp(\Delta\tau\,K)\,z_n,\qquad \Delta\tau=\frac{2\pi}{\mathrm{den}}=\frac{2\pi}{N}\]
and \(\exp(\Delta\tau K)\) is a real orthogonal matrix. Because the same
orthogonal matrix acts on the real and imaginary parts, the total norm
\(\|z\|^2\) and the square-closure \(z^{\mathsf T}z=\sum_e z_e^2\) are
exactly conserved.

\textbf{Verification} The error by which
\(\exp(\Delta\tau K(\varphi_n))z_n\) reproduces the canonical orbit
\(z_{n+1}\) is at most \(3\times10^{-15}\) (\(N=6,22,40\), \(\le300\)
steps), confirming that \(K\) is precisely the generator (Paper 1).

\subsubsection{2.2 The single-wave equation of
motion}\label{the-single-wave-equation-of-motion}

Writing the generator flow \(dz/d\tau=Kz\) in components and decomposing
\(z_e=r_e e^{i\varphi_e}\), from the real and imaginary parts of
\(\sum_f A_{ef}\sin(\varphi_f-\varphi_e)r_f e^{i(\varphi_f-\varphi_e)}\),
for adjacent edges \(f\sim e\) we obtain
\[\boxed{\ \frac{dr_e}{d\tau}=\frac12\sum_{f\sim e}r_f\,\sin\!\big(2(\varphi_f-\varphi_e)\big)\ }\tag{amplitude}\]
\[\boxed{\ r_e\,\frac{d\varphi_e}{d\tau}=\sum_{f\sim e}r_f\,\sin^2(\varphi_f-\varphi_e)\ }\tag{phase}\]
Each edge wave is a Kuramoto-type oscillator that couples only with
vertex-sharing adjacent edges via the sine of the phase difference.
Since \(\sin^2\ge0\), the phase advances monotonically (= rotates), and
the difference between waves gives the phase displacement modulo rigid
rotation. The increase/decrease of amplitude is driven by the second
harmonic \(\sin(2\Delta\varphi)\). In the frame modulo rigid rotation,
the amplitude rearrangement appears as reorientation (precession) of the
phase field, and this is the dynamical substance of the plane tilting
during the inflation era (§7).

\textbf{Verification} The amplitude/phase components of the small
substep \((\exp(\varepsilon K)z-z)/\varepsilon\) agree with the above
equations to at most \(10^{-6}\) (from rounding in the
\(\varepsilon\)-scale difference) (\(N=6,22,40\), Paper 1).

\begin{figure}
\centering
\pandocbounded{\includegraphics[keepaspectratio,alt={Figure 1: Verification of the single-wave equation of motion. The amplitude/phase components of the small substep (exp(εK)z−z)/ε (measured LHS) agree with the derived equations (RHS) on y=x (N=6).}]{figs/p0_capstone_en_fig1.png}}
\caption{Figure 1: Verification of the single-wave equation of motion.
The amplitude/phase components of the small substep (exp(εK)z−z)/ε
(measured LHS) agree with the derived equations (RHS) on y=x (N=6).}
\end{figure}

\textbf{Figure 1} Verification of the single-wave equation of motion
(N=6, details in Paper 1). Left: amplitude, right: phase. The
measurements (vertical axis) and the derived equations (horizontal axis)
lie on \(y=x\).

\subsection{3. Derivation of the high-symmetry floor------the imposed
conditions are the map plus the maximal-symmetry selection
criterion}\label{derivation-of-the-high-symmetry-floorthe-imposed-conditions-are-the-map-plus-the-maximal-symmetry-selection-criterion}

We show that the initial value (the floor) is not conveniently chosen
for the result, as a derivation that merely imposes the selection
criterion ``relative equilibrium and maximally symmetric'' on the map
(there is no other input). This claim is doubly backed by the
initial-value control experiments (Paper 2) and the algebraic exact
solution (Paper 3).

\subsubsection{3.1 The relative-equilibrium
condition}\label{the-relative-equilibrium-condition}

The floor is a rigidly rotating relative equilibrium, i.e., it satisfies
amplitude freezing \(dr_e/d\tau=0\) (\(\forall e\)) and phase velocity
uniform over all \(e\). From the amplitude equation, the former gives,
as a necessary condition,
\[\sum_{f\sim e}r_f\,\sin\!\big(2(\varphi_f-\varphi_e)\big)=0\qquad(\forall e)\tag{3.1}\]
(phase uniformity is automatically satisfied because, when the floor is
taken to be an eigenmode of \(iK\) in §3.3, \(z(t)=e^{-i\mu\tau}z_0\)).

\subsubsection{3.2 The 90-degree lattice is forced as the maximally
symmetric relative
equilibrium}\label{the-90-degree-lattice-is-forced-as-the-maximally-symmetric-relative-equilibrium}

There are two ways to realize (3.1): (i) each term vanishes individually
(termwise), (ii) individual terms are nonzero but the sum cancels
(cancellation). The \textbf{maximally symmetric equilibrium that
annihilates all coupling terms individually} requires
\(\sin(2\Delta\varphi)=0\), i.e., from
\[\sin(2\psi)=0\iff \psi=k\cdot90^\circ\ (k\in\mathbb Z)\] that all
adjacent phase differences are multiples of \(90^\circ\), that is, all
phases lie on a common 90-degree lattice \(\{0,90,180,270\}^\circ\)
(modulo global phase), which is necessary and sufficient. This is not a
choice but a consequence of the amplitude drive being the second
harmonic (§4, Paper 1).

\subsubsection{3.3 The amplitude is a deterministic eigenmode (algebraic
exact
solution)}\label{the-amplitude-is-a-deterministic-eigenmode-algebraic-exact-solution}

Given the 90-degree lattice, the generator is fixed to an integer
antisymmetric matrix (entries \(0,\pm1\)), and the floor \(z\) is
determined as its \(\sigma_{\max}\) eigenmode (the eigenvector of the
\textbf{most-negative eigenvalue in absolute value}
\(\mu=-\sigma_{\max}\) of \(iK\); the phase advances at
\(d\varphi/d\tau=-\mu=\sigma_{\max}>0\), consistent with the phase
equation of §2.2 \(r\dot\varphi=\sum r_f\sin^2\Delta\varphi\ge0\);
confirmed by real data and code). The amplitude is the square root of
the eigenvalues of the integer symmetric matrix \(K^{\mathsf T}K\) (=
the singular values of the integer antisymmetric matrix \(K\)) = an
algebraic number, a deterministic construction containing no continuous
tuning. \textbf{More generally, for
\(D=\operatorname{diag}(e^{i\varphi_e})\), on the 90-degree skeleton the
identity \(\sin\psi\,e^{i\psi}=i\sin^2\psi\) gives \(D^{-1}KD=iW\)
(\(W_{ef}=A_{ef}\sin^2\psi_{ef}\), real symmetric, nonnegative, 0/1), so
the complex self-consistent floor problem \(Kz=i\omega z\) (\(z=Dr\)) is
exactly equivalent to the eigenvalue problem \(Wr=\omega r\) of a real
symmetric nonnegative integer matrix. If \(W\) is irreducible, by
Perron--Frobenius the positive-amplitude principal eigenvector is unique
(\(\omega=\rho(W)=\sigma\))------that is, a self-consistent floor can be
constructed on the 90-degree skeleton, and its uniqueness is guaranteed
(Paper 1; the \(W\) of the 90-degree skeleton is the parity bipartite
graph \(W_{ef}=A_{ef}\mathbf 1_{p_e\neq p_f}\), and Paper 3 proves that
it is connected = irreducible for any \(N\ge3\)). The integer-K floor is
nothing but the closed form of the Perron eigenvector of this \(W\).
Zero-closure \(\sum z^2=0\) also follows automatically on the 90-degree
skeleton from \(\|x\|^2=\|y\|^2\) of this bipartite eigenmode (Paper 3,
for all nonzero eigenmodes, not only the Perron one).} This construction
has no free parameters, and \(\sigma^2\) (even \(N\): \(N(N-2)\) / odd
\(N\): \(N^2-2N-1\)) and each edge amplitude² (rational number, fixed by
parity class) can be written in closed form (only the parity
\(p_e=(i+j)\bmod2\) of the species \(k_e=(i+j)\bmod4\) makes
\(W_{ef}=A_{ef}\mathbf 1_{p_e\neq p_f}\) bipartite, derived from the
reduced Perron problem {[}even \(2\times2\) / odd \(3\times3\){]} for
\textbf{arbitrary \(N\)}; floor residual \(\sim10^{-15}\), sympy
symbolic + numerical \(N=3\)--\(40\) verification, Paper 3).
Furthermore, the active dimension is \(O(N)\) (of the \(M=O(N^2)\), the
nonzero spectrum is \(O(N)\) in count), and for \(N\ge6\) the closed
form \(\operatorname{rank}K=2(N-1)\), kernel dimension \((N-1)(N-4)/2\)
holds (\(N=3,4,5\) are low-dimensional exceptions: \(N=3\) is rank 2 /
kernel 1, \(N=4\) is rank 2 / kernel 4, \(N=5\) is rank 6 / kernel 4).
The edge assignment of the amplitude class and the odd-\(N\) residual
centroid, which were open in Paper 3, were also settled in closed form.
The make\_parent floor used by the experimental data (Papers 2 and 5--8)
is an unstable self-consistent eigenmode on the same 90-degree phase
lattice, but is a \textbf{different member of the relative-equilibrium
moduli} with different amplitude distribution and centroid (centroid
\(|\sum z|\): make\_parent \(0.25\)--\(1.41\), integer-K floor
\(\le0.24\)). The integer-K floor is not the exact version of the data
floor itself but an algebraically exact construction of a floor
satisfying the same ignition qualification (differences and corrections
in the section ``Differences from and corrections to cited papers'').

\begin{figure}
\centering
\pandocbounded{\includegraphics[keepaspectratio,alt={Figure 2b: The 90-degree theoretical floor (integer-K σ\_max analytic solution, all N=3..40). Phases 0/90/180/270, exact relative equilibrium (residual \textasciitilde1e-15), 4\textbar N gives exactly zero centroid.}]{figs/p0_capstone_en_fig2.png}}
\caption{Figure 2b: The 90-degree theoretical floor (integer-K σ\_max
analytic solution, all N=3..40). Phases 0/90/180/270, exact relative
equilibrium (residual \textasciitilde1e-15), 4\textbar N gives exactly
zero centroid.}
\end{figure}

\textbf{Figure 2b} Algebraic exact solution of the 90-degree floor
(\(N=3\)--\(40\), details in Paper 3). All waves lie on two orthogonal
axes, the 90-degree phase lattice.

\subsubsection{3.4 Zero-closure is not a requirement but a consequence
(Theorem)}\label{zero-closure-is-not-a-requirement-but-a-consequence-theorem}

\textbf{Theorem} If \(K\) is real antisymmetric and \(v=a+ib\) is a
self-consistent eigenmode \(iK(v)v=\mu v\) (\(\mu\neq0\)), then
\(\sum_e v_e^2=0\), that is, \(a^{\mathsf T}a=b^{\mathsf T}b\) and
\(a^{\mathsf T}b=0\).

\textbf{Proof} From the real and imaginary parts of
\(iK(a+ib)=iKa-Kb=\mu a+i\mu b\), \(Ka=\mu b,\ Kb=-\mu a\). By
antisymmetry,
\(\mu(a^{\mathsf T}a-b^{\mathsf T}b)=0\Rightarrow a^{\mathsf T}a=b^{\mathsf T}b\),
and
\(\mu a^{\mathsf T}b=a^{\mathsf T}Ka=0\Rightarrow a^{\mathsf T}b=0\). ∎

The proof uses only ``\(K\) is real antisymmetric'' and ``\(v\) is an
eigenvector of \(iK\) (\(\mu\neq0\))''; self-consistency is the
condition for \(v\) to be a floor and is not needed for the algebra.
\textbf{Neither zero-closure nor any particular amplitude distribution
is imposed on the construction of the floor. What is imposed is only the
map and the selection criterion `maximally symmetric relative
equilibrium (termwise + \(\sigma_{\max}\))', and the closure appears as
a theorem.} (Measured: \(\sum z^2\sim10^{-14}\) at the floor.
Factorization into zero-closure blocks is §3.5.)

\subsubsection{3.5 The factorization theorem of zero-closure and the
2-adic selection
rule}\label{the-factorization-theorem-of-zero-closure-and-the-2-adic-selection-rule}

\textbf{On the cyclic-Fourier high-symmetry floor}
(\(z_{ij}=r_d e^{i2\pi(i+j)/N}\)), zero-closure \(\sum z^2=0\) can be
exactly factorized into smaller zero-closure blocks (this is a
proposition about this high-symmetry floor and the root-of-unity
structure of the same distance class, not about the data floor in
general). The squared waves \(w=z^2\) within a distance class cycle
around a regular \(L\)-gon with phase step \(q=e^{i8\pi/N}\) (order
\(L=N/\gcd(N,4)\)), and the minimal zero-closure block is the smallest
prime factor of \(L\),

\[p=\operatorname{spf}(L),\qquad \boxed{\,p=2\iff L\ \text{even}\iff 8\mid N\iff \nu_2(N)\ge3\,}\]

(minimality by the vanishing-sum theorem of Lam--Leung). For \(N=8k\)
the minimal unit drops to two waves (\(\pm i\) photon pairs), and
otherwise it is a \(p\)-wave condensate (3 waves for \(3\mid N\), \(N\)
waves for odd prime \(N\)). That \(p=2\) holds only for \(8\mid N\) is
the \textbf{2-adic selection rule}: the square map doubles the phase
step \(4\pi/N\to8\pi/N\), and the condition for antipodal points (the
order-2 subgroup) to exist in \(\mathbb Z_L\) is \(L\) even
\(=\nu_2(N)\ge3\) (Paper 4). Importantly, closure \(\sum z^2=0\)
\textbf{alone} has a continuous moduli as solution space
(\(\mathcal C_k\simeq Q_{k-2}/S_k\), \(\dim_{\mathbb C}=k-2\) for
\(k\ge3\), Paper 9) and no discrete block size arises; it is quantized
to \(p=\operatorname{spf}(L)\) \textbf{only when the discrete symmetry
of the floor is imposed}. Note that the data floor make\_parent has no
exact \(\pm i\) pairs (only near pairs, see the section ``Differences
from and corrections to cited papers'')------the factorization depends
not on zero-closure or the 90-degree lattice but on the additional
structure of the cyclic-Fourier floor.

\subsubsection{3.6 No seed was planted------the ignition qualification
is an unstable relative
equilibrium}\label{no-seed-was-plantedthe-ignition-qualification-is-an-unstable-relative-equilibrium}

That the 90-degree floor is not a ``special configuration that happens
to work'' is established by the initial-value control experiments (Paper
2). Replacing only the initial value with three families and running the
same physics, (A) the symmetric parent with phases \(0/90\) (equal
amplitude) and (B) the zero-centroid parent both do not ignite (a
neutral fixed point, or immediate relaxation at step1); only (C) the
high-symmetry theoretical floor (cyclic-Fourier relative equilibrium)
and the experimental data's make\_parent floor ignite from the floor
with an exponential ramp. That is, the \textbf{ignition qualification is
``an unstable relative equilibrium (self-consistent eigenmode, residual
at machine precision)'' only}; zero centroid, square-closure, and high
symmetry themselves are not qualifications. The failure of (A)(B) is
established to be ``because it was \textbf{not an unstable relative
equilibrium},'' not ``because it was highly symmetric'' ((A) also
includes a neutral fixed point, and its not igniting demonstrates this
distinction). This is the initial-value control version of the
``amplification \(\iff\) unstable relative equilibrium'' of the
onset-mode paper {[}2, Concept 21798854{]}. In addition, the seed that
ignites the inflation is an infinitesimal, generic perturbation (whether
floor residual \(\sim10^{-13}\) or random \(10^{-8}\)), not a designed
one (§5 exhaustive experiment). \textbf{This point is established by an
experiment that directly verified the floor of the seed at 100-digit
precision} (precision/seed-isolation experiment, \(N=7,8\), \(D=N\)):
even generating the initial value at 100 digits and lowering it to
closure \(\sim10^{-102}\) and floor \(f_\perp\sim10^{-200}\), the
amplification rate \(\gamma_\tau\) \textbf{agrees to 8 digits (\(N=8\))
/ 6 digits (\(N=7\))} with the 64-bit run (= the exponential instability
is intrinsic to the dynamics and depends on neither precision nor seed
depth), it \textbf{necessarily ignites} after growing about 179 digits
from the floor \(\sim10^{-200}\) by a single exponential law (\(N=8\)
step 967 / \(N=7\) step 871), and the onset merely lags in proportion to
the number of seed digits
(\(\Delta\text{step}=\Delta(\text{seed digits})/(\gamma_{\rm step}/\ln10)\)).
The float64 rounding during evolution is only a secondary effect
advancing the onset by +6 steps. Therefore ignition is \textbf{neither a
double-precision artifact nor a planted seed}, and the floor is a
generically unstable starting point (the seed-independence of the
amplification rate is also consistent with the prior inflation paper
{[}Reference 1{]} and the seed-condition paper's \(\delta\) sweep
\(10^{-15}\)--\(10^{-2}\) {[}Concept DOI 10.5281/zenodo.21874481{]}).

\subsubsection{3.7 Established scope and open
matters}\label{established-scope-and-open-matters}

The following are open because they have not yet been derived from this
mathematical system (not out of timidity but because they are
underived). What determines the overall scale \(r^2\) of the floor is
not within the framework (scale is privately held). The normalization
\(\|z\|=1\) is a convention. Also, with a different generating measure a
different self-consistent fixed point exists (the moduli of relative
equilibria), and the 90-degree maximally symmetric floor is the most
symmetric member of that family, not the unique solution. \(N=3\) splits
in behavior by floor bifurcation: the T-type (\(\sigma=\sqrt2\),
\(|\sum z|=0.707\), ignites; make\_parent floor / integer-K floor) and
the triplet type (\(\mu/r^2=-3/2\), rank 1, neutral; cyclic-Fourier
floor). The \(N=3\) that runs in the data series is the T-type and
ignites (see the section ``Differences from and corrections to cited
papers''). These are explicitly noted as open.

\subsection{4. Why the self-consistent floor becomes 90-degree------the
reason make\_parent naturally became
90-degree}\label{why-the-self-consistent-floor-becomes-90-degreethe-reason-make_parent-naturally-became-90-degree}

We state the core of §3.2 independently. The amplitude drive is the
second harmonic \(\sin(2\Delta\varphi)\) arising from the
\(\sin(\varphi_f-\varphi_e)\) (first harmonic) of the generator, whose
zeros on \([0,360^\circ)\) are only \(0,90,180,270^\circ\) (confirmed by
enumeration). Therefore the maximally symmetric self-consistent relative
equilibrium that freezes the amplitude \textbf{termwise} (each term
individually) is limited to phase differences of all pairs being
multiples of \(90^\circ\)------a common 90-degree lattice. That
make\_parent converges for all \(N\) to a termwise 90-degree floor is an
\textbf{empirical fact} (Experiment 4-1 / Study 28). The equation of
motion in this section analytically explains \textbf{why the reached
floor is always a 90-degree lattice}: \textbf{the maximally symmetric
relative equilibrium that freezes the amplitude termwise is limited to
the 90-degree lattice by \(\sin2\Delta\varphi=0\) (Theorem)}. However,
the basin / the iteration-selection mechanism itself, by which the
make\_parent iteration selects this termwise branch rather than the
cancellation type, is a separate question (open). Paper 1 treats the
construction of the floor (self-consistent = eigenvalue problem of the
amplitude) with the 90-degree as \textbf{given} (an empirical fact that
naturally arose), and this capstone gives as a theorem the
\textbf{reason} why the reached floor is a 90-degree lattice (termwise ⇒
90-degree)------the empirical given and the analytic necessity are two
faces of the same fact.

\textbf{Measured (\(N=6\))} At the floor, the individual terms
\(|\sin(2\Delta\varphi)|\) of adjacent edges are at most
\(2.3\times10^{-12}\) (termwise equilibrium). By contrast, at lock (the
tilted vacuum), the individual terms are \(O(1)\), at most \(1.0\), yet
the weighted sum per edge is \(1.4\times10^{-7}\) (cancellation
equilibrium), and the phase differences deviate from the 90-degree
lattice by median \(23.5^\circ\) / max \(44.6^\circ\). That is, the
90-degree lattice is ``the unique maximally symmetric configuration
where the amplitude drive vanishes for every term individually,'' and
the locked vacuum is a cancellation-type equilibrium distinct from it
(Paper 1).

\begin{figure}
\centering
\pandocbounded{\includegraphics[keepaspectratio,alt={Figure 2: Left, the zeros of the amplitude drive sin(2Δφ) are limited to k·90°. Right, the measured distribution of the adjacent-edge individual terms \textbar sin2Δφ\textbar{} for the floor (90-degree lattice, individual terms concentrated at 0 = termwise) and lock (O(1), sum cancels = cancellation) (N=6).}]{figs/p0_capstone_en_fig3.png}}
\caption{Figure 2: Left, the zeros of the amplitude drive sin(2Δφ) are
limited to k·90°. Right, the measured distribution of the adjacent-edge
individual terms \textbar sin2Δφ\textbar{} for the floor (90-degree
lattice, individual terms concentrated at 0 = termwise) and lock (O(1),
sum cancels = cancellation) (N=6).}
\end{figure}

\textbf{Figure 2} 90-degree necessity (N=6, details in Paper 1).

\subsection{5. Spontaneous symmetry breaking and the degenerate
vacuum}\label{spontaneous-symmetry-breaking-and-the-degenerate-vacuum}

\subsubsection{5.1 Mechanism}\label{mechanism}

The floor that runs by SSB is a self-consistent unstable relative
equilibrium with a 90-degree skeleton (the data floor make\_parent, §3).
An infinitesimal seed \(\delta\) breaks the symmetry, is exponentially
amplified along the unstable eigendirection of the floor (inflation,
§6), and locks into a tilted relative equilibrium of equal amplitude and
phase spread. This is a Mexican-hat spontaneous symmetry breaking: floor
= an unstable maximum, valley = the family of degenerate relative
equilibria (not a finite discrete symmetry orbit, §5.2), the seed
spontaneously selecting one of its points. \textbf{What is called
``vacuum'' in this section is a mathematical designation for the
degenerate minima of the effective potential of the order parameter
\(r\) (= the locked relative equilibrium), not a characterization of the
physical vacuum (§9).}

From the measured dynamics of the order parameter \(r=\sqrt{H_\perp/H}\)
(\(H_\perp\) is the energy of the component orthogonal to the floor
plane), the Landau-type effective potential
\[V(r)=-\tfrac12\lambda_0 r^2+\frac{\lambda_0}{4\,r_{\rm vac}^2}r^4\]
can be reconstructed with measured parameters (\(\lambda_0=\ln\rho_0\),
\(\rho_0=\rho(J_0)\), \(r_{\rm vac}=\sqrt{H_\perp/H|_{\rm sat}}\))
(\(N=6\): \(\lambda_0=0.338,\ r_{\rm vac}=0.485\)). This is an effective
description (ansatz) with coefficients identified from measured
quantities. A first-principles derivation of \(V(r)\) is left as future
work (Paper 5).

\begin{figure}
\centering
\pandocbounded{\includegraphics[keepaspectratio,alt={Figure 3a: N=6. On the reconstructed effective potential V(r), the measured order parameter r=√(H⊥/H) falls from the floor (r=0, unstable summit) to the valley r\_vac=0.485 (color = step). On the sombrero surface, an unstable summit and a degenerate vacuum ring (5 seed points).}]{figs/p0_capstone_en_fig4.png}}
\caption{Figure 3a: N=6. On the reconstructed effective potential V(r),
the measured order parameter r=√(H⊥/H) falls from the floor (r=0,
unstable summit) to the valley r\_vac=0.485 (color = step). On the
sombrero surface, an unstable summit and a degenerate vacuum ring (5
seed points).}
\end{figure}

\begin{figure}
\centering
\pandocbounded{\includegraphics[keepaspectratio,alt={Figure 3b: N=40. Same procedure. λ0=0.076, r\_vac=0.330. Fall from floor to valley + degenerate vacuum ring.}]{figs/p0_capstone_en_fig5.png}}
\caption{Figure 3b: N=40. Same procedure. λ0=0.076, r\_vac=0.330. Fall
from floor to valley + degenerate vacuum ring.}
\end{figure}

\textbf{Figure 3} Mexican-hat SSB (measured reconstruction, details in
Paper 5). Top = N=6 (\(\lambda_0=0.338,\ r_{\rm vac}=0.485\)), bottom =
N=40 (\(\lambda_0=0.076,\ r_{\rm vac}=0.330\)).

\subsubsection{\texorpdfstring{5.2 Exhaustive verification
(\(N=6,\ldots,40\))}{5.2 Exhaustive verification (N=6,\textbackslash ldots,40)}}\label{exhaustive-verification-n6ldots40}

To the same floor \(Z_0\), five infinitesimal seeds of amplitude
\(10^{-8}\) in different orientations were applied, and each was run to
adaptive lock by the exact one\_step. In all 35 values × 5 seeds each =
\textbf{175 runs} for \(N=6,\ldots,40\): - Every seed causes inflation
(clear departure from the floor). - All lock into an isomorphic
equilibrium with equal amplitude
(\(|a_{\max}/a_{\min}-1|\lesssim10^{-3}\)) and a 1-step rotation angle
matching \(\Delta_{\rm tail}=-2\pi/N\) (error \(\sim10^{-5}\)). - The
phase configuration (vacuum orientation) differs by seed (median vacuum
difference \(4^\circ\)--\(48^\circ\)).

That is, from the unstable floor, an infinitesimal seed spontaneously
selects one of the degenerate vacua. This is the spontaneous symmetry
breaking of this system (\textbf{confirmed in 175/175}, Paper 5). The
\textbf{degenerate vacuum family} is \textbf{not identifiable at least
as a finite discrete orbit consisting of \(S_N\) + global phase +
complex conjugation}: it was confirmed that the five terminal states for
\(N=6\) do not map to one another by any of the \(6!=720\) vertex
permutations + global phase + complex conjugation (minimal residual
\(0.05\)--\(0.83\), Paper 5). \textbf{The point of this result is that
from the same 90-degree floor the terminal plane (vacuum orientation) is
not uniquely determined but is seed-dependent.} The
existence/identification of a larger continuous symmetry group, and the
continuity of the degenerate vacuum family (whether a 1-parameter
continuous scan of the seed direction
\(\delta(\theta)=\cos\theta\,\delta_1+\sin\theta\,\delta_2\) moves the
terminus continuously), are open (this open matter does not demote SSB
itself).

\begin{figure}
\centering
\pandocbounded{\includegraphics[keepaspectratio,alt={Figure 4: N=6, same floor + five different infinitesimal seeds. Left = the locked configuration of each seed (rigid rotation removed; equal radius but differing phase configuration = degenerate vacuum), center = all seeds match Δ=-2π/N, right = the N-dependence of the vacuum difference over all N=6-40.}]{figs/p0_capstone_en_fig6.png}}
\caption{Figure 4: N=6, same floor + five different infinitesimal seeds.
Left = the locked configuration of each seed (rigid rotation removed;
equal radius but differing phase configuration = degenerate vacuum),
center = all seeds match Δ=-2π/N, right = the N-dependence of the vacuum
difference over all N=6-40.}
\end{figure}

\textbf{Figure 4} SSB degenerate vacua (real data, N=6).

\begin{figure}
\centering
\pandocbounded{\includegraphics[keepaspectratio,alt={Figure 5: The SSB batch over all N=6-40. Left = the N-dependence of the vacuum difference (median/max), right = the measured Δ\_tail matches the theoretical -2π/N for all N (error \textasciitilde1e-5), amplitude ratio -1 \textasciitilde1e-3.}]{figs/p0_capstone_en_fig7.png}}
\caption{Figure 5: The SSB batch over all N=6-40. Left = the
N-dependence of the vacuum difference (median/max), right = the measured
Δ\_tail matches the theoretical -2π/N for all N (error
\textasciitilde1e-5), amplitude ratio -1 \textasciitilde1e-3.}
\end{figure}

\textbf{Figure 5} SSB over all \(N=6\)--\(40\) (175/175 confirmed).

\begin{figure}
\centering
\pandocbounded{\includegraphics[keepaspectratio,alt={Figure 5a: The inflation H⊥/H (step500) for all N=3-40 in the same arrangement as Figure 5b. The den=N (2π/N, canonical) curve is this system's inflation, ramping exponentially from floor \textasciitilde1e-29 → saturation. The others are denominator-control controls.}]{figs/p0_capstone_en_fig8.png}}
\caption{Figure 5a: The inflation H⊥/H (step500) for all N=3-40 in the
same arrangement as Figure 5b. The den=N (2π/N, canonical) curve is this
system's inflation, ramping exponentially from floor
\textasciitilde1e-29 → saturation. The others are denominator-control
controls.}
\end{figure}

\textbf{Figure 5a} Inflation \(H_\perp/H\) for all \(N=3\)--\(40\)
(step500, same arrangement as Figure 5b). In each panel the den\(=N\)
(\(2\pi/N\)) curve is this system's inflation, ramping exponentially
from the floor (\(\sim10^{-29}\), float64 rounding is the seed) →
saturation. As \(N\) increases, the rate drops and the onset lengthens
(§6). The other curves are denominator-control controls
(\(2\pi/(N\pm k)\), \(2\pi/124\)). From Figure 5a (inflation) → Figure
5b (terminal complex), floor → rapid expansion → lock can be surveyed
for all \(N\) at a glance.

\begin{figure}
\centering
\pandocbounded{\includegraphics[keepaspectratio,alt={Figure 5b: The terminus (step500) complex plane of the make\_parent data floor (all N=3..40). Equal-amplitude ring (radius \textasciitilde1/√M) with the phase spread = a locked tilted relative equilibrium. Small N is already locked, large N is still filling the ring (all lock by 2000 steps, §5.3).}]{figs/p0_capstone_en_fig9.png}}
\caption{Figure 5b: The terminus (step500) complex plane of the
make\_parent data floor (all N=3..40). Equal-amplitude ring (radius
\textasciitilde1/√M) with the phase spread = a locked tilted relative
equilibrium. Small N is already locked, large N is still filling the
ring (all lock by 2000 steps, §5.3).}
\end{figure}

\textbf{Figure 5b} Terminal complex diagram (make\_parent data floor,
step500, all \(N=3\)--\(40\)). The floor (Figure 2b of §3.3, and
make\_parent step0 = Figure 1 of Paper 2) has locked by SSB into a
tilted relative equilibrium of \textbf{equal amplitude} (radius
\(\sim1/\sqrt M\)) and \(\Delta=-2\pi/N\). Against the petal-like step0
of the 90-degree lattice, the terminus has each wave spread on a ring of
common radius (the \(a_\infty=\sqrt{H/M}\) of §8). Small \(N\) is
already locked, and large \(N\) (\(\ge18\)) is still filling the ring at
step500 and locks entirely by 2000 steps (§5.3). The integer-K analytic
floor (Figure 2b) also locks into the same universal terminus
(\(\Delta=-2\pi/N\), equal amplitude).

\subsubsection{\texorpdfstring{5.3 Long-time verification (\(N=22,40\),
2000
steps)}{5.3 Long-time verification (N=22,40, 2000 steps)}}\label{long-time-verification-n2240-2000-steps}

The large \(N\) (\(N=22,30,\ldots,40\)) that was unlocked at step500 was
due to truncation, and running without stopping to 2000 steps locks them
all. \(N=22\) fully locks at step \(663\), and \(N=40\) at step \(1299\)
(\(\Delta_{\rm tail}\) exactly \(-2\pi/N\), equal-amplitude ratio
\(1.0000\), floor/tail rigid-rotation residual machine-zero). A control
test confirmed that the initial value bit-matches the old run at step0
(step1 difference \(\sim10^{-16}\), thereafter inflation exponentially
amplifies), ensuring faithful reproduction (Paper 5).

The inflation curve (exponential growth of \(H_\perp/H\) → saturation)
and the terminal step2000 complex diagram are shown for \(N=22,40\).

\begin{figure}
\centering
\pandocbounded{\includegraphics[keepaspectratio,alt={Figure 5c: N=22 inflation H⊥/H (2000 steps). Grows by exponential ramp from the floor (\textasciitilde1e-29) and saturates.}]{figs/p0_capstone_en_fig10.png}}
\caption{Figure 5c: N=22 inflation H⊥/H (2000 steps). Grows by
exponential ramp from the floor (\textasciitilde1e-29) and saturates.}
\end{figure}

\textbf{Figure 5c} Inflation \(H_\perp/H\) (\(N=22\), 2000 steps). Grows
by exponential ramp from the floor (\(\sim10^{-29}\), float64 rounding
is the seed) and saturates.

\begin{figure}
\centering
\pandocbounded{\includegraphics[keepaspectratio,alt={Figure 5d: N=40 inflation H⊥/H (2000 steps). Similarly exponential ramp → saturation (saturation \textasciitilde3e-2).}]{figs/p0_capstone_en_fig11.png}}
\caption{Figure 5d: N=40 inflation H⊥/H (2000 steps). Similarly
exponential ramp → saturation (saturation \textasciitilde3e-2).}
\end{figure}

\textbf{Figure 5d} Inflation \(H_\perp/H\) (\(N=40\), 2000 steps).
Exponential ramp → saturation (as \(N\) increases the rate drops and the
onset lengthens, §6).

\begin{figure}
\centering
\pandocbounded{\includegraphics[keepaspectratio,alt={Figure 5e: N=22 terminal step2000 complex plane. Equal-amplitude ring (radius \textasciitilde1/√231) with the phase spread = fully locked.}]{figs/p0_capstone_en_fig12.png}}
\caption{Figure 5e: N=22 terminal step2000 complex plane.
Equal-amplitude ring (radius \textasciitilde1/√231) with the phase
spread = fully locked.}
\end{figure}

\textbf{Figure 5e} Terminal complex diagram (\(N=22\), step2000). All
\(M=231\) waves are spread on a ring of common radius
\(\sim1/\sqrt{231}\) = fully locked (step663).

\begin{figure}
\centering
\pandocbounded{\includegraphics[keepaspectratio,alt={Figure 5f: N=40 terminal step2000 complex plane. All 780 waves on a ring of radius \textasciitilde1/√780 = fully locked.}]{figs/p0_capstone_en_fig13.png}}
\caption{Figure 5f: N=40 terminal step2000 complex plane. All 780 waves
on a ring of radius \textasciitilde1/√780 = fully locked.}
\end{figure}

\textbf{Figure 5f} Terminal complex diagram (\(N=40\), step2000). All
\(M=780\) waves are spread on a ring of common radius
\(\sim1/\sqrt{780}=0.0358\) = fully locked (step1299,
\(\Delta_{\rm tail}=-2\pi/40\)). Figure 5b (large \(N\) at step500 still
filling the ring) locks all around by 2000 steps as in Figures 5e and
5f.

\subsection{6. The inflation curve = linear instability of the
floor}\label{the-inflation-curve-linear-instability-of-the-floor}

The inflation curve is the linear-instability amplification of an
infinitesimal seed. Write the spectral radius of the floor's linearized
map (Jacobian \(J\)) as \(\rho_0=\rho(J_0)\), \(\lambda_0=\ln\rho_0\) (a
quantity distinct from the \(iK\) eigenvalue \(\mu\) of §3; henceforth
the Jacobian side is consistently \(\rho_0,\lambda_0\)). \(\rho_0\)
gives
\[\frac{H_\perp}{H}(t)\approx H_{\perp0}\,\rho_0^{2t}\ \ (H_{\perp0}=\text{seed}^2),\qquad
\text{ramp slope}=2\log_{10}\rho_0,\qquad
\text{onset}\approx\frac{\ln(1/H_{\perp0})}{2\ln\rho_0}=\frac{\ln(1/\text{seed})}{\lambda_0}\]
The slope of \(H_\perp\) obtained by evolving a small \(\delta\) under
the fixed floor \(J\) matches \(2\log_{10}\rho_0\) exactly (\(N=6\):
\(0.2935\) vs \(0.2937\)). \(\rho_0\) decreases with \(N\) (\(N=6\):
\(1.402\), \(8\): \(1.269\), \(10\): \(1.185\)), so inflation becomes
slower and the onset lengthens (measured onsets are \(64,87,133\) steps
for \(N=6,8,10\)). A closed-form \(N\)-dependence law of the onset is
left as future work.

The amplification is a \textbf{two-stage dynamics}. Near the floor
\(\delta_{t+1}=J_0\delta_t\) with the initial/frozen rate
\(\lambda_0=\ln\rho_0\) (= the Landau summit curvature
\(V''(0)=-\lambda_0\), §5). Away from the floor the state itself changes
and \(\delta_t=\Phi(t,0)\delta_0\),
\(\Phi(t,0)=J(z_{t-1})\cdots J(z_0)\) (the \textbf{tangent cocycle} =
the time-ordered product of the time-varying, noncommuting Jacobians).
\textbf{The real inflation rate is the finite-time Lyapunov rate of this
tangent cocycle}

\[\gamma_t=\frac1t\ln\frac{\|P_\perp\Phi(t,0)\delta_0\|}{\|P_\perp\delta_0\|},\qquad \text{ramp slope}=\frac{2\gamma_t}{\ln10}\]

(\(P_\perp\) = the orthogonal projection onto the floor plane).
Reproduction A \textbf{directly confirmed} that the \(H_\perp\) slope
evolved by \(\prod_t J(z_t)\) along the real orbit matches the measured
ramp (time-varying / measured \(=0.999/1.004/1.019\), \(N=6,8,10\)).
\(J\) is non-normal (\(0.53\)--\(0.76\)) and noncommuting, so
\(\Phi(t,0)\neq J_0^t\) and the direction of maximal amplification is
amplified while rotating. \(\lambda_0=\ln\rho_0\) (the initial/frozen
rate near the floor) and the real rate \(\gamma_t\) are distinct
quantities, and what carries \(\gamma_t\) is \(\Phi(t,0)\) (the
amplification dynamics of the real orbit). For \(N=6,8,10\),
\(\gamma_t>\lambda_0\) (the real ramp exceeds the floor value by
\(18\)--\(31\%\))------not a ``correction'' but a consequence of the
two-stage dynamics. That \(\gamma_t\) does not in general fall below
\(\lambda_0\) (that \(\lambda_0\) is a lower bound) is limited in this
system to the observations for \(N=6,8,10\) and has no general proof. A
closed form of \(\gamma_t\) is also underived (Paper 6).

\begin{figure}
\centering
\pandocbounded{\includegraphics[keepaspectratio,alt={Figure 6: Inflation = linear instability. Left = the N-dependence of the floor-Jacobian ρ0=ρ(J0) and the onset (larger N decelerates, onset increases), center = the analytic ramp slope 2log10 ρ0 (floor) vs measured (non-normal gap), right = the analytic slope superimposed on the measured H⊥/H ramp for N=6.}]{figs/p0_capstone_en_fig14.png}}
\caption{Figure 6: Inflation = linear instability. Left = the
N-dependence of the floor-Jacobian ρ0=ρ(J0) and the onset (larger N
decelerates, onset increases), center = the analytic ramp slope 2log10
ρ0 (floor) vs measured (non-normal gap), right = the analytic slope
superimposed on the measured H⊥/H ramp for N=6.}
\end{figure}

\textbf{Figure 6} Inflation = linear instability of the floor (details
in Paper 6).

\subsection{7. Geometry after locking------plane switching on the
Grassmann and the 4-dimensional span of the initial and terminal
planes}\label{geometry-after-lockingplane-switching-on-the-grassmann-and-the-4-dimensional-span-of-the-initial-and-terminal-planes}

We establish what inflation and locking do in state
space------\textbf{the instantaneous rank-2 plane \(P(t)\) departs
exponentially from the floor plane on the Grassmann manifold
\(\mathrm{Gr}(2,M)\) and switches to a different rank-2 plane}. This is
the geometric consequence of the dynamics of §2 (amplitude rearrangement
= reorientation of the phase field = precession).

\textbf{(i) Instantaneous rank 2 and Grassmann motion} The map is a real
orthogonal rotation \(z_t=M_t z_0\) (\(Q_t=\exp(\Delta\tau K_t)\) is the
one-step map, \(M_t=Q_{t-1}\cdots Q_0\) the cumulative map, both real
orthogonal acting identically on real and imaginary parts), and the
state stays on a 2-dimensional plane (rank 2) at every instant (the
singular values of \([\Re z,\Im z]\) are only two-valued at all steps;
by conservation of \(z^{\mathsf T}z=0\) an orthogonal isometric frame).
Modulo global phase, the observable is the 2-dimensional subspace
\(P(t)=\mathrm{span}\{\Re z,\Im z\}\) itself (motion on
\(\mathrm{Gr}(2,M)\)), and the projection \(\Pi_t\) moves exactly under
the one-step map \(Q_t\) as \(\Pi_{t+1}=Q_t\Pi_tQ_t^{\mathsf T}\). At
the floor \(Q_0P_0=P_0\) (\(Ka=\mu b,\ Kb=-\mu a\)) is fixed (principal
angles zero), during inflation \(\Pi_t\neq\Pi_0\) moves, and at
saturation it re-locks to \(P_{\rm term}\).

\textbf{(ii) The initial and terminal planes span four dimensions} The
locked plane \(P_{\rm term}\) tilts from \(P_0\), and the principal
angles are, independently of \(N\), \textbf{both nonzero} (\(N=6\):
\(27.1^\circ,30.8^\circ\), \(N=40\): \(18.1^\circ,20.4^\circ\)), with
\(\dim(P_0+P_{\rm term})=4\). The tilt of the locked plane from the
floor is \(14.2^\circ\)--\(29.0^\circ\) (median \(\sim19.2^\circ\)) over
all \(N=6\)--\(40\). \textbf{From the same floor this terminal
orientation is not uniquely determined but is seed-dependent} (§5).

\textbf{(iii) The tilt, the two principal angles, and \(H_\perp/H\) are
the same Grassmann quantity (an identity holding at all times)} By
definition, for the principal angles \(\theta_1(t),\theta_2(t)\) between
\(P_0\) and \(P(t)\) at each time \(t\),

\[\frac{H_\perp(t)}{H}=\frac12\left(\sin^2\theta_1(t)+\sin^2\theta_2(t)\right)=\frac14\big\|\Pi_t-\Pi_0\big\|_F^2\]

(\(\Pi=UU^{\mathsf T}\)) holds (not only for the terminal
\(\Pi_{\rm term}\) but for all \(t\)). Hence from the initial ramp
\(H_\perp/H\propto e^{2\gamma_t t}\) of §6, immediately

\[\boxed{\ \|\Pi_t-\Pi_0\|_F\propto e^{\gamma_t t}\ }\]

------\textbf{inflation is exponential departure from the floor plane on
\(\mathrm{Gr}(2,M)\)}. This is not ``it can be read that way'' but a
\textbf{derived} geometric description combining the identity from the
definition with the exponential rate of §6.

\textbf{(iv) The 4D span is a minimal span, not \(SO(4)\) dynamics}
During the transition the second direction rises (transient
\(s_3/s_1\approx0.10\)--\(0.20\) = the local span swept in a short time
by the two moving planes), and after lock it returns to a single plane
(tail \(s_3/s_1\approx0\)). The cumulative map
\(M_t=\prod\exp(\Delta\tau K)\) does not preserve the 4-dimensional
subspace (\(M_tS\neq S\), confirmed by a discriminating computation), so
\textbf{``the dynamics is an \(SO(4)\) double rotation on four
dimensions'' is rejected}------\(\dim(P_0+P_{\rm term})=4\) is the
minimal span needed for the relative configuration of two 2-planes, not
an invariant space in which the dynamics lives. A short-window SVD test
separated out that the plane tilt is real, not an artifact of
long-window aggregation (Paper 7).

Note that when initially the complex plane is only a single sheet, its
normal is in a degenerate orthogonal complement and is unobservable
(conceptually). Only when inflation raises a specific orientation in the
normal space does the relative phase (tilt) between the original plane
and the new plane become readable. What can be read is always a
relational quantity (the relative phase between planes); the absolute
normal direction cannot be recovered from the state alone (\(Q_t\) hides
the orientation, and no conserved quantity points to it).

\subsection{8. Terminal universality, selective erasure of memory, and
anonymity}\label{terminal-universality-selective-erasure-of-memory-and-anonymity}

\textbf{Terminal universality (intensive state quantity)}: after
locking, the terminus is of equal amplitude, and its radius is
determined only by the extensive variables \((H,M)\):
\(a_\infty=\sqrt{H/M}\). This is an exact consequence of
\(H=\sum_e|z_e|^2=Ma_\infty^2\); \((H,M)\) are extensive quantities,
\(\rho=H/M\) is intensive, and \(a_\infty=\sqrt\rho\) is an
\textbf{intensive terminal state quantity}. What is nontrivial is not
the formula but the \textbf{terminal universality}------a mixed system
with the triangular-number mixture \(M(9)=M(7)+M(6)\) normalized to the
same density, and a single system, both converge to the same
\(a_\infty=1/6\); and even though the presence/absence of ignition
itself changes by initial-value family (staged / analytic floor / v2)
(v2 is a non-igniting neutral fixed point), the terminus is common
(Paper 8).

\textbf{Selective erasure of memory (Axis B)}: the transient (inflation
era) records the initial structure in the floor Jacobian \(J_0\) and the
tangent cocycle \(\Phi(t,0)=\prod_t J(z_t)\) and determines the onset
and path. But after nonlinear saturation \(|z_e|^2\to H/M\), and
\textbf{the configuration information disappears from the amplitude
sector}. On the other hand, the SSB history remains in the phase
configuration and the terminal plane orientation (§5, §7). Therefore

\[\boxed{\text{the amplitude universalizes to }H/M\text{ and forgets the configuration, while the phase configuration and plane orientation retain the SSB history}}\]

------not ``forgetting everything'' but a \textbf{sector-selective
erasure} (selective erasure of history information).

\textbf{Anonymity and relationalism (Axis C)}: the edge numbering is a
gauge of the vertex permutation \(S_N\), and only \(S_N\)-invariants can
be read (``which wave'' cannot be read). At the terminus even the
individuality of \(|z_e|\) disappears, so they cannot be distinguished
even by value, and the indistinguishability of identical particles
appears as a consequence of the dynamics. This is the same principle as
the relational readout of §7 (the absolute normal is unobservable, only
the relative phase between planes can be read), and \textbf{the physical
content resides only in relational invariants modulo \(S_N\), global
phase, and coordinate rotation} (Paper 9).

\textbf{Closure does not generate discrete species------the quantization
is carried by the floor}: even under the \textbf{hypothesis} that the
minimal closure block, not the individual wave, is the unit of
information, closure \(\sum z^2=0\) \textbf{alone} does not generate
discrete species (a mathematical fact):
\(\mathcal C_k\simeq Q_{k-2}/S_k\) (the projective quadric hypersurface
\(\dim_{\mathbb C}=k-2\)), where \(k=2\) has \([1{:}i],[1{:}{-}i]\)
collapsing into one orbit under \(S_2\) (\(\pm\) is not an anonymous
bit), and \(k\ge3\) is continuous. \(Q_{k-2}\) is connected for
\(k\ge3\), and the quotient by the finite group \(S_k\) is also
connected, so from closure alone \textbf{no ``species'' as a discrete
connected component arises}. What is discrete is the block size
\(p=\operatorname{spf}(L)\) of §3.5------\textbf{closure + the discrete
symmetry of the cyclic-Fourier high-symmetry floor quantizes the block
size} (Paper 9). The identification of the self-consistent dynamics
(orientation, winding, topology, dynamical selection rule) that selects
a discrete type from the continuous relational moduli is the next
research problem.

\subsection{9. Established scope and open
problems}\label{established-scope-and-open-problems}

The following are open because they have not yet been derived from this
mathematical system (not out of timidity but because they are
underived).

\begin{itemize}
\tightlist
\item
  \textbf{Identification of coordinates}: whether the 4-dimensional span
  obtained in §7 can be identified with physical \(xyzt\) (or
  \(xyztRQ\), \(xyztRQ_1Q_2Q_3\)) is open. What is established is that
  the dimension of the minimal span spanned by the initial 2 planes and
  the terminal 2 planes is \(\dim(P_0+P_{\rm term})=4\), and that their
  relative configuration has two nonzero principal angles
  \(\theta_1,\theta_2\neq0\)------not four physical axes, and not two
  independent rotational dynamics (\(SO(4)\) dynamics is rejected in
  §7).
\item
  \textbf{Gauge (graduation)}: whether a gauge can be defined along each
  direction, and its type (equal-interval / logarithmic / winding), is
  undefined. The external step \(\tau\) is not a physical internal clock
  (prior work measured that the equal-spacing of \(\tau\) breaks after
  inflation), and whether an intrinsic equal-interval clock exists is
  open. A clock stands only as the beat of a coherent,
  distinct-frequency two-component; it does not stand for a single
  coherent mode.
\item
  \textbf{Narrowness of the state's symmetry}: this research primarily
  targets a self-consistent single-coherent relative equilibrium with a
  90-degree skeleton and its small anisotropic perturbations. States
  with coexisting particles of various masses and charges are not
  treated.
\item
  \textbf{Vacuum}: whether the 90-degree floor is the true vacuum or a
  condensate that already has structure is unspecified. The vacuum
  itself is not characterized.
\item
  \textbf{Caveat on circular derivation}: the terminal equal amplitude
  (Paper 8) inherently involves this engine's phase normalization (the
  processing that kills the amplitude), so it is not presented as an
  independent corroboration of the norm form. Independent verification
  in an amplitude-preserving control dynamics is open.
\item
  \textbf{The band/hair register of Paper A is scaffolding (demotion of
  the assessment)}: the types (boson/fermion, etc.) read in the prior
  work (Paper A, Concept 21874481), which gives each wave
  \(16\times8=128\) internal registers, are scaffolding by
  \textbf{declared constants} not derived from the closure axiom or
  anonymity, and this series does not cite them as a physical claim. To
  physicalize the types, one must derive a discrete degree of freedom
  from the closure axiom + the discrete symmetry of the floor + a
  dynamical selection rule (currently not achieved, Paper 9).
\item
  \textbf{Quantization mechanism of discrete species}: the
  self-consistent dynamics (orientation, winding, topology, dynamical
  selection rule) that selects a discrete type corresponding to a
  particle species from the continuous closure moduli
  (\(\mathcal C_k\simeq Q_{k-2}/S_k\)) is unidentified (§8, Paper 9).
\item
  \textbf{Others}: the privately-held scale \(r^2\), the moduli of
  relative equilibria, the \(N=3\) exception, and the non-readability of
  the absolute time \(t\).
\end{itemize}

\subsection{10. Related work (structural agreement------correspondence
to prior art, not the derivation
source)}\label{related-work-structural-agreementcorrespondence-to-prior-art-not-the-derivation-source}

The results of this research were derived independently of the map of
§2, but their structure agrees with existing theory.

\begin{itemize}
\tightlist
\item
  \textbf{Kuramoto model} (Kuramoto 1975): the generator
  \(K_{ef}=A_{ef}\sin(\varphi_f-\varphi_e)\) is a Kuramoto-type phase
  coupling on \(L(K_N)\), and locking corresponds to synchronization.
\item
  \textbf{Spontaneous symmetry breaking / Landau theory /
  Goldstone-Higgs} (Landau 1937; Goldstone 1961; Anderson 1963; Higgs
  1964): the phenomenon of §5 is an instance of SSB on a discrete map,
  and the reconstructed effective potential is a Ginzburg--Landau-type
  sombrero.
\item
  \textbf{Transient amplification of non-normal matrices /
  pseudospectra} (Trefethen \& Embree 2005): structural agreement of the
  transient amplification of a fixed non-normal system. In this system,
  the amount by which the real ramp persistently exceeds the floor
  eigenvalue is carried not by the transient of the fixed floor but by
  the time-varying, noncommuting tangent cocycle \(\Phi(t,0)\) along the
  real orbit, as confirmed in Reproduction A (§6).
\item
  \textbf{Spectra of line graphs} (Cvetković et al., algebraic graph
  theory): the eigenvalues of \(A=L(K_N)\).
\item
  \textbf{Relative equilibria} (Marsden \& Ratiu, geometric mechanics):
  the rigidly rotating relative equilibria of the floor and lock.
\end{itemize}

\subsection{Differences from and corrections to cited papers (notes for
consistency)}\label{differences-from-and-corrections-to-cited-papers-notes-for-consistency}

There are three discrepancies between this capstone and the foundational
papers regarding the identification of the floor. So that the claims of
the capstone do not waver, we explicitly state the claim on the
original-paper side, and that it is erroneous (or how it should be read
with restriction), with the reason from derivation/measurement. The
description of the original papers at the time of writing is left
without weakening, and errors are corrected here as errors. What is
established is that the ignition qualification is ``an unstable relative
equilibrium'' only, and that the floors satisfying it split into three
mutually distinct families: (i) the \textbf{make\_parent floor} used by
the experimental data (Papers 5--8) (90-degree phase skeleton, unstable
relative equilibrium, nonzero centroid, no exact \(\pm i\) pairs), (ii)
the \textbf{integer-K analytic floor} of Paper 3 (an exact relative
equilibrium on the same 90-degree phase skeleton, centroid \(\le0.24\),
with \(\pm i\) pairs for \(N=8k\)), (iii) the \textbf{cyclic-Fourier
floor} of Paper 2(C) / Paper 4 (a different high-symmetry
relative-equilibrium family, generally not a 90-degree phase lattice but
based on cyclic symmetry------\(z_{ij}=r_d e^{i2\pi(i+j)/N}\)). (i)(ii)
are distinct members sharing the 90-degree skeleton, and (iii) is a
third family whose skeleton itself differs.

\textbf{Correction 1: the implication of Paper 3 that ``the integer-K
floor is the exact version of the make\_parent floor (the same floor)''
turned out to be erroneous.} Paper 3, in its Abstract, §1, and §4
(precision comparison table), treated the make\_parent floor as ``a
solution approximating the self-consistent eigenmode by iteration
(residual \(\sim10^{-13}\)),'' and, placing the integer-K floor's
relative-equilibrium residual \(\sim10^{-15}\) as ``exceeding the
\(\sim10^{-13}\) of make\_parent,'' presented the integer-K floor as
\textbf{a more precise construction of the same floor}. This implication
of ``the same floor'' \textbf{turned out to be erroneous} by the
following measurements: (i) the centroid \(|\sum z|\) differs between
the two floors (make\_parent \(0.25\)--\(1.41\) / integer-K floor
\(\le0.24\), with the integer-K floor \(\sim10^{-15}\) for \(4\mid N\)).
(ii) the \(\pm i\) pair structure differs (Paper 4 / Experiment 15:
make\_parent has no exact \(\pm i\) pairs / the integer-K and
cyclic-Fourier floors have them for \(N=8k\)). Correctly, the two floors
are not approximate/exact versions of the same floor but
\textbf{distinct members} of the relative-equilibrium moduli (the family
of self-consistent fixed points selected by the generating measure)
satisfying the same ignition qualification (unstable relative
equilibrium, 90-degree phase lattice). The individual facts of Paper 3
(the make\_parent floor is an approximate relative equilibrium, the
integer-K floor is an exact relative equilibrium, the magnitude of the
residuals) are correct, but the implication of regarding the two floors
as the same floor was erroneous.

\textbf{Correction 2: the \(\pm i\) photon-pair factorization is a
property of the high-symmetry theoretical floor (cyclic-Fourier floor)
and does not hold exactly for the data floor.} Difference: the full
\(\pm i\)-pair coverage (\(N=8k\)) and the \(p\)-wave condensate of
Paper 4 are properties of the cyclic-Fourier floor
\(z_{ij}=r_d e^{i2\pi(i+j)/N}\). The make\_parent data floor has no
exact \(\pm i\) pairs (Paper 4 / Experiment 15: strict pairs 0 for all
\(N\le40\), best \(\varepsilon=3.71\times10^{-8}\)). Correction: read
the factorization of capstone §3.5 with the floor family
restricted------exact for the high-symmetry theoretical floor, only near
pairs for the data floor.

\textbf{Correction 3: \(N=3\) splits into ignition/neutral by floor
bifurcation.} Difference: \(N=3\) has two bifurcations, the T-type
(\(\sigma=\sqrt2\), \(|\sum z|=0.707\), ignites; make\_parent floor /
integer-K floor) and the triplet type (\(\mu/r^2=-3/2\), rank 1,
neutral; cyclic-Fourier floor). Correction: the ``\(N=3\) is the
\(-3/2\) exception'' of capstone §3.7 refers to the triplet-type
bifurcation. The \(N=3\) that runs in the data series is the T-type and
ignites at onset 45 (Paper 2 / Experiment 13).

By the above, the capstone's claims about the floor are unified to
``data floor = make\_parent floor (90-degree lattice, unstable relative
equilibrium, nonzero centroid, no exact \(\pm i\) pairs).'' The
integer-K analytic floor of Paper 3 is positioned as a distinct member
on the same 90-degree skeleton (not an exact version but a different
floor of the same qualification), and the cyclic-Fourier floor of Paper
4 as a different high-symmetry relative-equilibrium family not sharing
the 90-degree skeleton. The igniting branches of the three share the
ignition qualification ``an unstable relative equilibrium,'' but are
distinguished by the floor's skeleton, centroid, and \(\pm i\) pair
structure. However, the cyclic-Fourier triplet type of \(N=3\)
(\(\mu/r^2=-3/2\), rank 1) is a neutral branch and does not ignite
(Correction 3).

\subsection{References}\label{references}

\begin{enumerate}
\def\labelenumi{\arabic{enumi}.}
\tightlist
\item
  Noriaki Kihara, ``Mechanism of inflationary rapid expansion in
  self-consistent relational-wave closed systems,'' Version DOI
  10.5281/zenodo.22176949 / Concept DOI 10.5281/zenodo.22112008. (The
  source for the definition of the dynamics of this paper, the
  kinematization of zero-closure, and \(\tau\neq\) time.)
\item
  This series, Papers 1--9 (Foundations of Self-Consistent
  Relational-Wave Closed Systems, Concept DOIs in Appendix A).
\item
  Y. Kuramoto, \emph{Chemical Oscillations, Waves, and Turbulence},
  Springer (1984); ``Self-entrainment of a population of coupled
  non-linear oscillators'', 1975.
\item
  L. D. Landau, ``On the theory of phase transitions'', 1937; J.
  Goldstone, Nuovo Cimento 19 (1961) 154; P. W. Anderson, Phys. Rev.~130
  (1963) 439; P. W. Higgs, Phys. Rev.~Lett. 13 (1964) 508.
\item
  L. N. Trefethen and M. Embree, \emph{Spectra and Pseudospectra},
  Princeton (2005).
\item
  D. Cvetković, P. Rowlinson, S. Simić, \emph{An Introduction to the
  Theory of Graph Spectra}, Cambridge (2010).
\item
  J. E. Marsden and T. S. Ratiu, \emph{Introduction to Mechanics and
  Symmetry}, Springer (1999).
\end{enumerate}

\subsection{Appendix A: List of self-citation
papers}\label{appendix-a-list-of-self-citation-papers}

Because of the volume, this is consolidated in an appendix rather than
the main text.

\subsubsection{(a) Foundational papers of this series (Papers 1--9; the
sources for the details, exhaustive data, and figures of each section of
this
capstone)}\label{a-foundational-papers-of-this-series-papers-19-the-sources-for-the-details-exhaustive-data-and-figures-of-each-section-of-this-capstone}

{\def\LTcaptype{none} % do not increment counter
\begin{longtable}[]{@{}
  >{\raggedright\arraybackslash}p{(\linewidth - 6\tabcolsep) * \real{0.2500}}
  >{\raggedright\arraybackslash}p{(\linewidth - 6\tabcolsep) * \real{0.2500}}
  >{\raggedright\arraybackslash}p{(\linewidth - 6\tabcolsep) * \real{0.2500}}
  >{\raggedright\arraybackslash}p{(\linewidth - 6\tabcolsep) * \real{0.2500}}@{}}
\toprule\noalign{}
\begin{minipage}[b]{\linewidth}\raggedright
Paper
\end{minipage} & \begin{minipage}[b]{\linewidth}\raggedright
Title (abbreviated)
\end{minipage} & \begin{minipage}[b]{\linewidth}\raggedright
Corresponding section
\end{minipage} & \begin{minipage}[b]{\linewidth}\raggedright
Version/Concept DOI
\end{minipage} \\
\midrule\noalign{}
\endhead
\bottomrule\noalign{}
\endlastfoot
Paper 1 & Single-wave equation of motion; condition for the
self-consistent floor on the 90-degree skeleton (amplitude eigenvalue
problem) & §2, §4 & Version 10.5281/zenodo.22728761 / Concept
10.5281/zenodo.22728760 \\
Paper 2 & Qualification screening of the initial value = unstable
relative equilibrium & §3.6 & Version 10.5281/zenodo.22728824 / Concept
10.5281/zenodo.22728823 \\
Paper 3 & Analytic exact solution of the 90-degree floor (integer-K
σmax, parity bipartite, arbitrary-N closed form, O(N)) & §3.3 & Version
10.5281/zenodo.22728996 / Concept 10.5281/zenodo.22728995 \\
Paper 4 & Zero-closure factorization theorem (L=N/gcd(N,4), p=spf(L),
2-adic selection rule) & §3.5 & Version 10.5281/zenodo.22729081 /
Concept 10.5281/zenodo.22729080 \\
Paper 5 & Spontaneous symmetry breaking and the degenerate vacuum (all
N6-40, 175/175, degenerate family not a finite discrete orbit) & §5 &
Version 10.5281/zenodo.22729089 / Concept 10.5281/zenodo.22729088 \\
Paper 6 & Inflation is the linear instability of the floor (finite-time
Lyapunov rate of the tangent cocycle) & §6 & Version
10.5281/zenodo.22729096 / Concept 10.5281/zenodo.22729095 \\
Paper 7 & Geometry after locking (Grassmann plane switching,
4-dimensional span, 2 nonzero principal angles, SO(4) rejected) & §7 &
Version 10.5281/zenodo.22729108 / Concept 10.5281/zenodo.22729107 \\
Paper 8 & Additivity law a∞=√(H/M) (intensive state quantity, selective
erasure of memory) & §8 & Version 10.5281/zenodo.22729116 / Concept
10.5281/zenodo.22729115 \\
Paper 9 & Anonymity and the minimal unit of information (study paper,
Cₖ≃Q\_\{k-2\}/S\_k) & §8 & Version 10.5281/zenodo.22729128 / Concept
10.5281/zenodo.22729127 \\
\end{longtable}
}

\subsubsection{(b) Already-published parent papers (the underlay of this
series)}\label{b-already-published-parent-papers-the-underlay-of-this-series}

{\def\LTcaptype{none} % do not increment counter
\begin{longtable}[]{@{}
  >{\raggedright\arraybackslash}p{(\linewidth - 4\tabcolsep) * \real{0.3333}}
  >{\raggedright\arraybackslash}p{(\linewidth - 4\tabcolsep) * \real{0.3333}}
  >{\raggedright\arraybackslash}p{(\linewidth - 4\tabcolsep) * \real{0.3333}}@{}}
\toprule\noalign{}
\begin{minipage}[b]{\linewidth}\raggedright
Paper
\end{minipage} & \begin{minipage}[b]{\linewidth}\raggedright
Version DOI
\end{minipage} & \begin{minipage}[b]{\linewidth}\raggedright
Concept DOI
\end{minipage} \\
\midrule\noalign{}
\endhead
\bottomrule\noalign{}
\endlastfoot
Mechanism of inflationary rapid expansion in self-consistent
relational-wave closed systems (v2) & 10.5281/zenodo.22176949 &
10.5281/zenodo.22112008 \\
Onset-mode discrimination (amplification ⟺ unstable relative
equilibrium, all 86 states) & --- & 10.5281/zenodo.21798854 \\
Thermodynamic readout from a self-consistent complex relational system
without assuming spacetime (Λ, cosmological term) &
10.5281/zenodo.22240034 & 10.5281/zenodo.22240033 \\
Readout of spacetime and wave metric from anonymous complex zero-closure
(log spacetime, c=1) & 10.5281/zenodo.22282218 &
10.5281/zenodo.22282217 \\
Conditions under which a seed generates particles and the lower bound of
resolution (Paper A) & 10.5281/zenodo.21874482 &
10.5281/zenodo.21874481 \\
Reproduction specification (capstone + 3 chapters) &
10.5281/zenodo.22317636 & (to be confirmed) \\
\end{longtable}
}

\subsubsection{(c) Not yet published but planned for mutual citation (to
be cross-cited upon publication); in-series
experiments}\label{c-not-yet-published-but-planned-for-mutual-citation-to-be-cross-cited-upon-publication-in-series-experiments}

\begin{itemize}
\tightlist
\item
  \textbf{Precision / seed-isolation experiment (100 digits)}
  \texttt{N7\_N8\_DeltaTau\_denominator\_wide\_sweep\_20260902/precision\_seed\_dynamics\_isolation\_100digit\_20260902/}
  (the source for the seed-floor verification of §3.6. At IC100, closure
  \(\sim10^{-102}\), floor \(f_\perp\sim10^{-200}\), \(\gamma_\tau\)
  agrees to 8 digits with IC64, grows about 179 digits from the floor by
  a single exponential law and fires. README / ANALYSIS /
  \texttt{results/precision\_summary.csv}). Related: multiprecision
  50-digit
  \texttt{Nall\_linear1000000\_steps5\_mpmath50\_equimodular\_selfconsistent\_directHperp\_20260828/}.
\item
  \textbf{Reproduction A: ramp excess = time-varying non-normal product
  (direct Lyapunov)}
  \texttt{位相ロック全N確認\_相対平衡\_20260912/ランプ直接Lyapunov\_非正規時変積\_20260912/}
  (direct confirmation of the non-normal correction mechanism of §6. The
  \(H_\perp\) slope of the time-varying product \(\prod J(z_t)\) matches
  the measurement and exceeds the floor freeze by 18--31\%,
  \(N=6,8,10\). \texttt{analyze\_ramp\_lyapunov\_20260912.py} /
  \texttt{results/ramp\_lyapunov\_summary.csv} / REPORT / SHA /
  run\_all).
\item
  Periodic table of waves\_completed paper\_v1 (charge clock / mass
  clock; planned source for the open clock theory of §9).
\item
  Discovery ledger \texttt{発見台帳\_干渉保存力学\_20260905起点.md} (the
  unified record of experiments and studies = the backbone of
  reproduction).
\end{itemize}

\subsection{Appendix B: List of figures, generating programs, and
relative paths
(reproducibility)}\label{appendix-b-list-of-figures-generating-programs-and-relative-paths-reproducibility}

We list the content of each figure, its \textbf{generating program}, and
its \textbf{figure relative path} (relative to this capstone file). The
\texttt{../} in the paths points to directly under this series
(\texttt{干渉保存力学\_資格審査とシード無し系列\_20260831/}). The
dynamics map is identical to the canonical
\texttt{run\_N3\_N40\_stage123\_v1.py} (SHA256 \texttt{1abf2353…}) in
all figures, with no physical change.

{\def\LTcaptype{none} % do not increment counter
\begin{longtable}[]{@{}
  >{\raggedright\arraybackslash}p{(\linewidth - 6\tabcolsep) * \real{0.2500}}
  >{\raggedright\arraybackslash}p{(\linewidth - 6\tabcolsep) * \real{0.2500}}
  >{\raggedright\arraybackslash}p{(\linewidth - 6\tabcolsep) * \real{0.2500}}
  >{\raggedright\arraybackslash}p{(\linewidth - 6\tabcolsep) * \real{0.2500}}@{}}
\toprule\noalign{}
\begin{minipage}[b]{\linewidth}\raggedright
Figure
\end{minipage} & \begin{minipage}[b]{\linewidth}\raggedright
Content
\end{minipage} & \begin{minipage}[b]{\linewidth}\raggedright
Generating program (relative path)
\end{minipage} & \begin{minipage}[b]{\linewidth}\raggedright
Figure relative path
\end{minipage} \\
\midrule\noalign{}
\endhead
\bottomrule\noalign{}
\endlastfoot
Figure 1 & Single-wave EOM verification (measured vs derived, \(y=x\)) &
\texttt{../位相ロック全N確認\_相対平衡\_20260912/単一波運動方程式\_導出と検証\_20260912/make\_figures\_eom\_90deg\_20260912.py}
(verification \texttt{verify\_single\_wave\_eom\_20260912.py}) &
\texttt{../位相ロック全N確認\_相対平衡\_20260912/単一波運動方程式\_導出と検証\_20260912/fig\_single\_wave\_eom\_verify.png} \\
Figure 2 & 90-degree necessity (\(\sin2\Delta\varphi\) zeros +
floor/lock distribution) & same as above
\texttt{make\_figures\_eom\_90deg\_20260912.py} (verification
\texttt{verify\_90deg\_necessity\_20260912.py}) &
\texttt{../位相ロック全N確認\_相対平衡\_20260912/単一波運動方程式\_導出と検証\_20260912/fig\_90deg\_necessity.png} \\
Figure 2b & Algebraic exact solution of the 90-degree floor (integer-K
floor, all-N complex diagram) &
\texttt{../90度理論床\_整数K解析解\_全N\_20260909/make\_90deg\_floor\_analytic\_v1.py}
(closed-form verification
\texttt{closed\_form\_90deg\_floor\_verify\_v1.py}, parity
\texttt{parity二部構造\_解析的証明\_20260912/verify\_parity\_bipartite\_analytic\_20260912.py})
&
\texttt{../90度理論床\_整数K解析解\_全N\_20260909/fig\_90deg\_floor\_complex\_planes\_analytic.png} \\
Figure 3 & Mexican-hat effective potential (N=6 / N=40) &
\texttt{../位相ロック全N確認\_相対平衡\_20260912/自発的対称性の破れ\_シード依存縮退真空\_20260912/make\_figure\_mexican\_hat\_20260912.py}
(N40 is \texttt{make\_figure\_mexican\_hat\_N40\_20260912.py}) &
\texttt{../位相ロック全N確認\_相対平衡\_20260912/自発的対称性の破れ\_シード依存縮退真空\_20260912/fig\_mexican\_hat\_ssb.png}
(N40: \texttt{fig\_mexican\_hat\_ssb\_N40.png}) \\
Figure 4 & SSB degenerate vacua (\(N=6\), real data) &
\texttt{../位相ロック全N確認\_相対平衡\_20260912/自発的対称性の破れ\_シード依存縮退真空\_20260912/make\_figure\_ssb\_20260912.py}
(verification \texttt{verify\_ssb\_seed\_vacua\_20260912.py}) &
\texttt{../位相ロック全N確認\_相対平衡\_20260912/自発的対称性の破れ\_シード依存縮退真空\_20260912/fig\_ssb\_degenerate\_vacua.png} \\
Figure 5 & SSB for all \(N\) 175/175 + \(\Delta=-2\pi/N\) &
\texttt{../位相ロック全N確認\_相対平衡\_20260912/自発的対称性の破れ\_シード依存縮退真空\_20260912/全N6\_40\_SSBバッチ\_20260912/make\_figure\_batch\_20260912.py}
(run \texttt{run\_ssb\_batch\_N6\_N40\_20260912.py}) &
\texttt{../位相ロック全N確認\_相対平衡\_20260912/自発的対称性の破れ\_シード依存縮退真空\_20260912/全N6\_40\_SSBバッチ\_20260912/fig\_ssb\_batch\_allN.png} \\
Figure 5a & All-N inflation \(H_\perp/H\) (step500, same arrangement as
complex diagram) &
\texttt{../make\_parent型初期値\_自己無撞着構造\_20260907/wrapper\_run\_makeparent\_sweep\_v1.py}
(run) + \texttt{postprocess\_makeparent\_sweep\_v1.py} (plotting) &
\texttt{../make\_parent型初期値\_自己無撞着構造\_20260907/full\_N3\_N40\_sweep/fig\_Hperp\_makeparent\_N3\_N40\_500.png} \\
Figure 5b & Terminal complex diagram (make\_parent floor step500, all
\(N\), equal-amplitude ring) &
\texttt{../make\_parent型初期値\_自己無撞着構造\_20260907/wrapper\_run\_makeparent\_sweep\_v1.py}
(run) + \texttt{plot\_complex\_planes\_makeparent\_v1.py} (plotting) &
\texttt{../make\_parent型初期値\_自己無撞着構造\_20260907/full\_N3\_N40\_sweep/fig\_complex\_plane\_step500\_makeparent\_N3\_N40.png} \\
Figure 5c & Inflation \(H_\perp/H\) (\(N=22\), 2000step) &
\texttt{../位相ロック全N確認\_相対平衡\_20260912/N22\_2000step検証走行\_20260912/plot\_inflation\_N22\_v1.py}
(run \texttt{wrapper\_run\_N22\_2000\_v1.py}) &
\texttt{../位相ロック全N確認\_相対平衡\_20260912/N22\_2000step検証走行\_20260912/fig\_Hperp\_N22\_den22\_2000.png} \\
Figure 5d & Inflation \(H_\perp/H\) (\(N=40\), 2000step) &
\texttt{../位相ロック全N確認\_相対平衡\_20260912/N40\_2000step検証走行\_20260912/plot\_inflation\_N40\_v1.py}
(run \texttt{wrapper\_run\_N40\_2000\_v1.py}) &
\texttt{../位相ロック全N確認\_相対平衡\_20260912/N40\_2000step検証走行\_20260912/fig\_Hperp\_N40\_den40\_2000.png} \\
Figure 5e & Terminal complex step2000 (\(N=22\), equal-amplitude ring) &
\texttt{../位相ロック全N確認\_相対平衡\_20260912/N22\_2000step検証走行\_20260912/plot\_complex\_plane\_final\_N22\_v1.py}
&
\texttt{../位相ロック全N確認\_相対平衡\_20260912/N22\_2000step検証走行\_20260912/fig\_complex\_plane\_step2000\_N22.png} \\
Figure 5f & Terminal complex step2000 (\(N=40\), equal-amplitude ring) &
\texttt{../位相ロック全N確認\_相対平衡\_20260912/N40\_2000step検証走行\_20260912/plot\_complex\_plane\_final\_N40\_v1.py}
&
\texttt{../位相ロック全N確認\_相対平衡\_20260912/N40\_2000step検証走行\_20260912/fig\_complex\_plane\_step2000\_N40.png} \\
Figure 6 & Inflation = linear instability (\(N\)-dependence, ramp,
superposition) &
\texttt{../位相ロック全N確認\_相対平衡\_20260912/インフレ曲線\_線形不安定解析\_20260912/make\_figure\_inflation\_20260912.py}
(analysis \texttt{analyze\_inflation\_linear\_instability\_20260912.py})
&
\texttt{../位相ロック全N確認\_相対平衡\_20260912/インフレ曲線\_線形不安定解析\_20260912/fig\_inflation\_linear\_instability.png} \\
\end{longtable}
}

\textbf{Reproducibility upon publication (mandatory)}: when publishing
this capstone, upload in bulk all generating programs of the above
table, the input data (floor npz, states npz), the output figures, each
package's \texttt{README} / \texttt{SHA256SUMS} / \texttt{run\_all.sh},
and the canonical \texttt{run\_N3\_N40\_stage123\_v1.py}, and keep the
state in which \texttt{run\_all.sh} can regenerate the figures and data.
The packages of foundational Papers 1--9 (Appendix A) must likewise be
complete.

\subsection{Peer-review reflection
log}\label{peer-review-reflection-log}

\begin{itemize}
\tightlist
\item
  v1→v1.1 (2026-09-12): correction of the coefficient of the §6 onset
  formula
  (\(\ln(1/H_{\perp0})/(2\ln|\mu|)=\ln(1/\text{seed})/\ln|\mu|\)); ``the
  unique termwise equilibrium'' → ``maximally symmetric equilibrium;
  global uniqueness not claimed''; §5 ``vacuum'' explicitly stated as a
  mathematical designation for the degenerate minima of the effective
  potential; \(V(r)\) explicitly noted as an ansatz; §7 separated into
  the three layers of instantaneous rank 2 / 4D span / double rotation;
  a caveat on designations added at the opening.
\item
  v1.1→(this series capstone, 2026-09-12): correction of the thinness of
  the content. With the direction and claims unchanged, §1 was thickened
  as a research program and causal chain, §3 integrated the
  qualification screening (Paper 2) and the algebraic exact solution
  (Paper 3, Figure 2b), §3.4 the factorization (Paper 4), and §8 added
  the caveat on circular derivation. The details, exhaustive data, and
  figure sources of each section were linked to foundational Papers
  1--9, and Appendix A (list of self-citations) was newly established.
\item
  Content review (2026-09-12): inspected consistency with the
  foundational papers and stated the three discrepancies in the
  identification of the floor explicitly in the section ``Differences
  from and corrections to cited papers.'' In particular, \textbf{the
  implication of Paper 3 that ``the integer-K floor = the exact version
  of the make\_parent floor (the same floor)'' was corrected as
  erroneous by the measurement of the centroid and \(\pm i\) pair
  structure} (the original papers' description left without weakening,
  errors stated explicitly as errors). Also removed ambiguous words,
  excuses, and poorly-grounded numbers (onset
  ``\(\approx7.4\)--\(7.9N\)'', etc.), and explicitly noted
  unestablished derivations (onset closed form, \(V(r)\) first
  principles, \(S^3\) identification, etc.) as future work.
\item
  Fixing of the sign convention (2026-09-12, cross-checked against real
  data / canonical code): the floor is the eigenmode of the
  \textbf{most-negative eigenvalue in absolute value}
  \(\mu=-\sigma_{\max}\) of \(iK\), with
  \(d\varphi/d\tau=-\mu=\sigma_{\max}>0\) (consistent with the phase
  equation \(r\dot\varphi=\sum r_f\sin^2\Delta\varphi\ge0\); for \(N=6\)
  the make\_parent floor \(\mu=-4.868\) / integer-K floor
  \(\mu=-4.899\), both the minimal eigenvalue of \(iK\)). The notation
  ``the maximal eigenvalue of \(iK\)'' of §3.3 and Paper 3 was
  sign-inconsistent and was corrected (the designation
  \(\sigma_{\max}=|\mu|\) = spectral radius = maximal singular value is
  retained; the form of the floor is \(z=\sigma_{\max}u-iKu\)). Note
  that the \(|\mu_{\max}|\) of §6 and Paper 6 is the spectral radius of
  the floor Jacobian \(J\) and is a quantity distinct from this \(\mu\).
\item
  SO(4) discriminating computation (2026-09-12, real data / canonical
  code): constructed the cumulative map
  \(M_t=\prod\exp(\Delta\tau K(\varphi_n))\) and inspected the
  invariance of the 4-dimensional subspace
  \(S=\mathrm{span}(P_0\cup P_{\rm term})\). \(M_tS=S\) breaks
  (out-of-\(S\) component \(0.84\)--\(0.99\), \(M_S\) non-orthogonal
  \(\|M_S^{\mathsf T}M_S-I\|=1.21\), \(N=6\)). Hence \textbf{``the
  dynamics is an \(SO(4)\) double rotation on four dimensions'' is
  rejected}, and §7 / Paper 7 were corrected to ``4-dimensional span, 2
  nonzero principal angles, double rotation of the relative
  configuration (the cumulative map is not \(SO(4)\)).'' The
  already-established \(\dim(P_0+P_{\rm term})=4\) and 2 nonzero
  principal angles are retained.
\item
  Reinforcement of the citation for the seed floor (2026-09-12): §3.5
  added \textbf{an experiment that verified the floor of the seed at
  100-digit precision}
  (precision\_seed\_dynamics\_isolation\_100digit\_20260902, \(N=7,8\)).
  It has been shown that even lowering to closure \(\sim10^{-102}\),
  floor \(\sim10^{-200}\) at IC100, \(\gamma_\tau\) agrees to 8 digits
  with 64-bit and necessarily fires. By this, the previous version's
  description ``verification at multiprecision is additional research''
  is erroneous, and was replaced with a citation as \textbf{verified}
  (removing the defensive hedge).
\item
  Re-review of open matters by cross-checking with existing papers
  (2026-09-12): cross-checked all ``open / additional research'' items
  against the discovery ledger and all packages. \textbf{Corrected 2
  items that had been treated as open despite being already
  verified}------(i) \(N\ge18\) lock (established by Paper 5 \(N=22,40\)
  + N=64 phase-lock verification, §Paper 2), (ii) transition rank-4
  \(S^3\) (inspected as transient with no persistent structure by
  short-window SVD + discriminating computation, Paper 7). \textbf{3
  items settled by reproduction}------(A) the non-normal correction
  mechanism = time-varying non-normal product (the slope of the
  time-varying product matches the measurement and exceeds the floor
  freeze by 18--31\%, asserted, Paper 6), (B) the terminal block
  decomposition does not stand for the data series (make\_parent floor →
  terminus) and the minimal block is intrinsic to the cyclic-Fourier /
  integer-K high-symmetry floor (rejected, Paper 9), (C) \(\sigma=N-1\)
  ↔ vertex sector does not correspond, being a coincidence of dimension
  labels (the 3-decomposition and \(\sigma\to N-1\) retained, rejected,
  Paper 8). Each reproduction was saved in a dedicated package (py +
  results + README + REPORT + SHA + run\_all) and cited in the relevant
  paper. Those left genuinely open (\(V(r)\) first principles, the
  closed form of the onset/amplification rate, coordinates/gauge/clock,
  the amplitude-preserving control dynamics for circular derivation) are
  explicitly noted as future work. The \(N\ge18\) long-time run of the
  theoretical floor was judged unnecessary as it is trivial from the
  \(N=64\) confirmation of the data floor.
\item
  Synchronization accompanying the arbitrary-\(N\) analytification of
  Paper 3 (2026-09-12): confirmed that only the parity
  \(p_e=(i+j)\bmod2\) of the species \(k_e=(i+j)\bmod4\) makes the \(W\)
  of Paper 1 the bipartite graph
  \(W_{ef}=A_{ef}\mathbf 1_{p_e\neq p_f}\) (the \(W\) of the iterated
  floor matches the \(W\) of species parity at machine precision,
  ‖difference‖ \(\le3\times10^{-27}\); the floor can be constructed in
  one shot by \(k\to W\to\) Perron \(\to z=Dr\) with no iteration, and
  matches the iterated floor modulo sign gauge). By this, the closed
  forms of Paper 3 (\(\sigma^2\), amplitude², full spectrum) were
  promoted to \textbf{arbitrary \(N\)} (sympy symbolic + numerical
  \(N=3\)--\(40\)) via the reduced Perron problem, and the active
  dimension \(O(N)\) (for \(N\ge6\), rank\(K=2(N-1)\), kernel
  \((N-1)(N-4)/2\); \(N=3,4,5\) low-dimensional exceptions rank 2 / 1,
  rank 2 / 4, rank 6 / 4), and the closed forms of the amplitude-class
  edge assignment (= parity class) and odd-\(N\) residual centroid,
  which were open in Paper 3, were settled. Furthermore, it was
  established as a general proof that zero-closure \(\sum z^2=0\) is a
  consequence of the bipartite eigenmode (\(\|x\|^2=\|y\|^2\), for all
  nonzero modes not only the Perron one), that the connectivity of the
  bipartite \(W\) holds for any \(N\ge3\) (recovering the Perron
  applicability condition of Paper 1), and that the sign-gauge family is
  \(2^{M-1}\). The Paper 3 citation of §3.3 was updated from
  ``\(N=3\)--\(20\) sympy'' to ``arbitrary \(N\),'' and the \(W\)
  irreducibility from ``\(N=3\)--\(40\) confirmed'' to ``proved for any
  \(N\ge3\).'' Reproduction
  \texttt{90度理論床\_整数K解析解\_全N\_20260909/parity二部構造\_解析的証明\_20260912/}.
\item
  \textbf{Reconstruction of the capstone after full review of 1--9
  (2026-09-13)}: integrated into the capstone the results of
  self-reviewing and promoting Papers 1--9 in order. The Abstract and §1
  were fully rewritten with the new causal chain (map → 90-degree
  relative equilibrium → ignition qualification → \textbf{inflationary
  rapid expansion} → SSB / plane reorientation → terminal
  universalization → relational invariants) + the central proposition +
  the unifying axes A/B/C (the old ``first\ldots{} second\ldots{}''
  redundant development deleted). §3.5 promoted the factorization
  theorem (\(p=\operatorname{spf}(L)\),
  \(p=2\iff8\mid N\iff\nu_2(N)\ge3\)) to an independent section, and
  demoted the ignition qualification to §3.6 and the open matters to
  §3.7 with renaming (``honest open matters'' → ``established scope and
  open matters''). §5's ``unstable maximally symmetric equilibrium'' →
  ``unstable relative equilibrium of the 90-degree skeleton
  (make\_parent),'' ``35/35'' → ``175/175,'' and reflected the vacuum
  manifold = continuous moduli (720-vertex-permutation test). §6 was
  changed from ``floor eigenvalue + non-normal correction'' to a
  \textbf{two-stage dynamics} (the finite-time Lyapunov rate of the
  tangent cocycle \(\Phi(t,0)=\prod_t J(z_t)\); \(\mu_{\max}\) is a
  lower bound). §7 was changed to the Grassmann formulation
  (\(\Pi_{t+1}=Q_t\Pi_tQ_t^{\mathsf T}\), the identity
  \(H_\perp/H=\tfrac14\|\Pi_t-\Pi_0\|_F^2\), exponential departure on
  \(\mathrm{Gr}(2,M)\)), and the title updated. A new §8 ``Terminal
  universality, selective erasure of memory, and anonymity'' was added
  (Papers 8 and 9 moved from auxiliary to body, Axes B/C). The old §8
  ``What is not claimed'' was renamed to §9 ``Established scope and open
  problems,'' with the 128 registers of Paper A explicitly demoted as
  scaffolding, and §9 related work → §10. The audit criterion was
  unified to ``what is proved was proved / the exhaustive was confirmed
  / the rejected was rejected / only the underived is open'' (removal of
  evasive language, \textbf{not weakening the phenomenon name of
  inflation by replacing it with a mechanism word}). Also added, from
  existing figures, the all-N inflation figure (Figure 5a, step500 grid)
  and the 2000-step inflation + terminal complex for N=22/40 (Figures
  5c--5f). Appendix B was expanded into a reproducibility table of
  ``figures, generating programs, relative paths'' and the bulk-upload
  policy upon publication was stated.
\item
  \textbf{Refinement review of the capstone (2026-09-13, 7 points.
  Separate theorem from measurement / restrict the subject)}: ①
  restricted the subject of the factorization of §3.5 / the Abstract to
  ``on the cyclic-Fourier high-symmetry floor'' (not the data floor in
  general). ② the reason for the failure of ignition in §3.6 from ``was
  not a relative equilibrium'' → ``was \textbf{not an unstable relative
  equilibrium}'' (including the neutral fixed point of (A)). ③ refined
  the SSB vacuum manifold's ``continuous moduli'' to ``\textbf{not
  identifiable as a finite discrete symmetry orbit (\(S_N\) + global
  phase + conjugation)},'' explicitly noting the continuity as open by a
  seed 1-parameter scan, and deleted ``\(G/H\) analytic identification
  is generally impossible'' (for a continuous group, \(G/H\) is a
  continuous orbit). The continuity of the C\_k abstract closure space
  (\(Q_{k-2}\)) is retained as it is proved. ④ separated §6's
  ``\(\mu_{\max}\) is a lower bound of the real rate'' into
  ``\(\lambda_0=\ln\rho_0\) (initial/frozen rate) and the real rate
  \(\gamma_t\) are distinct quantities; for \(N=6,8,10\),
  \(\gamma_t>\lambda_0\); no general proof.'' Unified the Jacobian-side
  notation to \(\rho_0=\rho(J_0),\lambda_0\) (resolving the clash with
  the \(\mu\) of \(iK\) in the text). ⑤ stated §7's identity
  \(H_\perp(t)/H=\tfrac14\|\Pi_t-\Pi_0\|_F^2\) at all times and
  strengthened \(\|\Pi_t-\Pi_0\|_F\propto e^{\gamma_t t}\) as derived. ⑥
  §9 ``maximally symmetric single-coherent mode'' → ``self-consistent
  single-coherent relative equilibrium of the 90-degree skeleton,'' and
  §10 Trefethen--Embree updated to the time-varying tangent-cocycle
  criterion. ⑦ synchronized old section-number references (factorization
  §3.4→§3.5, \(N=3\) exception §3.6→§3.7, 100-digit appendix §3.5→§3.6).
  The title was updated to the current position (local instability →
  inflation → SSB → Grassmann → terminal universality → relationalism).
\item
  \textbf{Whole-consistency audit just before commit (2026-09-13, 2
  substantive points + 1 notational point)}: ① revived the \(N=3,4,5\)
  exceptions in the rank formula of §3.3 / the review log------the
  active dimension \(O(N)\) holds in general, but the closed form
  \(\operatorname{rank}K=2(N-1)\), kernel \((N-1)(N-4)/2\) is limited to
  \(N\ge6\) (explicitly noting the low-dimensional exceptions \(N=3\):
  rank 2 / kernel 1, \(N=4\): rank 2 / kernel 4, \(N=5\): rank 6 /
  kernel 4). ② corrected the places that had mixed the cyclic-Fourier
  floor as ``another member of the 90-degree floor'' into the
  three-family separation------(i) make\_parent floor (90-degree
  skeleton), (ii) integer-K analytic floor (the same 90-degree
  skeleton), (iii) cyclic-Fourier floor (a different high-symmetry
  relative-equilibrium family, generally not a 90-degree lattice but
  cyclic symmetry \(z_{ij}=r_de^{i2\pi(i+j)/N}\)). Made the three
  families explicit at the opening and end of ``Differences from and
  corrections to cited papers,'' and restricted the subject of the
  block-size quantization of §1(c) / §8 to ``the discrete symmetry of
  the cyclic-Fourier high-symmetry floor.'' The three share the ignition
  qualification ``an unstable relative equilibrium'' but are
  distinguished by the floor skeleton, centroid, and \(\pm i\) pair
  structure. ③ unified the coefficient of the body of \(V(r)\) in §5.1
  to \(\lambda\to\lambda_0\) (\(\lambda_0=\ln\rho_0\), resolving the
  notational clash between the parameter definition and the body).
\end{itemize}

\subsection{Reproduction}\label{reproduction}

The verification scripts, exhaustive data, and additional figures for
each claim of this capstone are stored in a set (py + REPORT +
SHA256SUMS + run\_all.sh) in the packages of the corresponding
foundational papers (Papers 1--9) (relative paths in Appendix A /
Appendix B). The dynamics map is identical to the prior paper (Reference
1) (canonical \texttt{run\_N3\_N40\_stage123\_v1.py}, SHA256
\texttt{1abf2353fee2e4f56f05e7a6f149fd086885136beb61ab571b48a56b09691567}
cross-checked, physical equations unchanged). The series table of
contents is
\texttt{関係波閉鎖系基礎シリーズ\_論文一覧と元資料\_20260912.md}.

\end{document}
